\exercisegroup

\begin{enumerate}
\def\labelenumi{\arabic{enumi})}
\tightlist
\item
  设 \(S\) 是任意一个无穷集合。
\end{enumerate}

\begin{enumerate}
\def\labelenumi{\alph{enumi})}
\tightlist
\item
  证明：赋范空间 \(\mathcal{B}(\mathrm{S})\)（由 S 到 \(\mathbf{R}\) 中的有界映射组成；I, 第 4 页, 例）中任一完全集的最小基数等于 \(2^{\mathrm{Card}(\mathrm{S})}\)（考察 S 的诸子集的特征函数所成的集，并注意 \(\mathbf{R}\) 中存在一个处处稠密的可数集）。b) 证明：赋范空间 \(\ell^1 (\mathrm{S})\)（I, 第 4 页, 例）中任一完全集的最小基数等于 \(\mathrm{Card}(\mathrm{S})\)。
\end{enumerate}

\begin{enumerate}
\def\labelenumi{\arabic{enumi})}
\setcounter{enumi}{1}
\item
  在域 R 上的乘积拓扑向量空间 \(E = R^{N}\) 中，用 \(e_{n} (n \in \mathbf{N})\) 表示直和 \(\mathbf{R}^{(\mathbf{N})}\) 的典范基的诸元素。置 \(a_{0} = e_{0}\)，\(a_{n} = e_{0} + (1/n)e_{n}\)（对 \(n \geqslant 1\)）。证明：对每个整数 \(n \geqslant 0\)，诸 \(a_{i}\)（其中 \(0 \leqslant i \leqslant n\)）在 E 中构成一个拓扑无关族，但无穷族 \((a_{n})_{n \geqslant 0}\) 不是拓扑无关的。若 M 是闭向量子空间 \(Ra_{0}\)，则 \(\dot{a}_{n}\)（\(a_{n}\) 在 E/M 中的类，对 \(n \geqslant 1\)）构成一个拓扑无关族，但由诸 \(a_{n}\)（指标 \(n \geqslant 1\)）在 E 中生成的闭向量子空间 N 包含 M。
\item
  设 E 是 R 上的一个拓扑向量空间，f 是 E 的加法群到 R 中的一个同态。证明：若存在 E 中 0 的一个邻域，使 f 在其中有界，则 f 是 E 中的连续线性型。特别地，当 f 半连续（下半连续或上半连续）时即属此种情形。
\item
  用 K 表示赋有绝对值 \(p(\xi) = |\xi|^{1/2}\) 的域 R。设 E 是 K 上的向量空间，由定义在 I = \([0, 1]\) 上、处处右连续且在点 1 处为零的实值受控函数组成；证明：映射 \(x \mapsto \|x\| = \int_{0}^{1} |x(t)|^{1/2} dt\) 在 E 上是一个范数。证明：对 E 中每个函数 \(x \geqslant 0\)，在 E 中存在两个函数 \(x_{1} \geqslant 0\)、\(x_{2} \geqslant 0\)，使得 \(x = \frac{1}{2}(x_{1} + x_{2})\) 且 \(\|x_{1}\| = \|x_{2}\| = \frac{1}{\sqrt{2}} \|x\|\)。由此推出，E 上每个连续线性型都恒等于零。
\item
  设 K 是一个 Hausdorff 拓扑除环，其拓扑是局部逆有界的 (GT, III, § 6, 习题 22)。把 I, 第 12 页 的 prop. 2 与 I, 第 13 页 的 th. 1 推广到 K 上的拓扑向量空间；当 K 还是完备的时候，同样把 I, 第 13 页 的 th. 2 与 I, 第 14 页 的 prop. 3 加以推广。
\item
  设 K 是把 \(Q^{2}\) 的通常拓扑转移到域 \(\mathbf{Q}(\sqrt{2})\) 上（经由映射 \((x, y) \mapsto x + y\sqrt{2}\)）而得到的拓扑除环。
\end{enumerate}

\begin{enumerate}
\def\labelenumi{\alph{enumi})}
\item
  设 \(\mathbf{E}\) 是集合 \(\mathbf{Q}(\sqrt{2})\)，带其 \(\mathbf{K}\) 上的向量空间结构以及由 \(\mathbf{R}\) 的拓扑诱导的拓扑。证明 \(\mathbf{E}\) 是一个 Hausdorff 拓扑向量空间，在 \(\mathbf{K}\) 上的维数为 1，但它不同构于 \(\mathbf{K}_s\)。
\item
  设 \(\mathbf{F}\) 是拓扑向量空间 \(\mathbf{E} \times \mathbf{E}\)（在 \(\mathbf{K}\) 上）；在 \(\mathbf{F}\) 中，超平面 \(\mathbf{E} \times \{0\}\) 是闭的，但不存在连续线性型 \(f\)（在 \(\mathbf{E} \times \mathbf{E}\) 上），使该超平面由方程 \(f(x) = 0\) 给出。
\end{enumerate}

\begin{enumerate}
\def\labelenumi{\arabic{enumi})}
\setcounter{enumi}{6}
\item
  设 K 是一个非离散且非完备的赋值除环，E 是拓扑向量子空间 \(K + Ka\)（\(\hat{K}\) 中，其中 \(a \notin K\)），F 是乘积空间 \(K \times E\)。在 F 中，子空间 \(M = K \times \{0\}\) 是闭的且余维数为 2。设 N 是 F 中由向量 (0, 1) 与 (1, a) 生成的 M 的补子空间；证明 F 不是 M 与 N 的拓扑直和。
\item
  设 \(p\) 是一个素数，\(\mathbf{Q}_p\) 是 \(p\) 进数域 (GT, III, § 6, 习题 23)。设 \(\mathrm{E}_0\) 是拓扑乘积空间 \(\mathbf{Q}_p \times \mathbf{R}\)；若用 \(\mathbf{K}\) 表示带离散拓扑的域 \(\mathbf{Q}\)，则 \(\mathrm{E}_0\) 是 \(\mathbf{K}\) 上的拓扑向量空间。设 \(\mathbf{M}\) 是由元素 \((r, r)\)（其中 \(r\) 取遍 \(\mathbf{Q}\)）组成的向量子空间；又设 \(\theta\) 是一个无理数，\(\mathbf{N}\) 是由元素 \((0, r\theta)\)（其中 \(r\) 取遍 \(\mathbf{Q}\)）组成的向量子空间。设 \(\mathbf{E}\) 是子空间 \(\mathbf{M} + \mathbf{N}\)（在 \(\mathrm{E}_0\) 中）；证明 \(\mathbf{N}\) 是 \(\mathbf{E}\) 中的一个闭超平面，但不存在 \(\mathbf{N}\) 的拓扑补子空间（注意 \(\mathbf{M}\) 在 \(\mathrm{E}_0\) 中处处稠密）。
\item
  设 \(\mathbf{X}\) 是一个 Hausdorff 拓扑空间，\(\mathbf{V}\) 是一个维数 \(n\) 有限的向量子空间，含于空间 \(\mathcal{C}(\mathbf{X};\mathbf{R})\) 内。
\end{enumerate}

\begin{enumerate}
\def\labelenumi{\alph{enumi})}
\item
  证明：存在 \(n\) 个两两不相交的开集 \(\mathrm{U}_i\)（\(1 \leqslant i \leqslant n\)），位于 \(\mathbf{X}\) 中，使得若有函数 \(f \in \mathbf{V}\) 在每个 \(\mathrm{U}_i\) 内都恒等于零，则它在 \(\mathbf{X}\) 中也恒等于零（利用 A, II, § 7.5, cor. 3）。
\item
  设 \(x_{i} \in \mathbf{U}_{i}\)（\(1 \leqslant i \leqslant n\)）。由 a) 推出：存在一个常数 \(c > 0\)，使得对每个函数 \(f \in \mathbf{V}\)，有
\end{enumerate}

\[
\sup _ {x \in \mathrm{X}} | f (x) | \leqslant c \sum_ {i = 1} ^ {n} | f (x _ {i}) |.
\]

\begin{enumerate}
\def\labelenumi{\arabic{enumi})}
\setcounter{enumi}{9}
\tightlist
\item
  设 \(\mathbf{K}\) 是一个局部紧的非离散赋值除环，\(\mathbf{E}\) 是 \(\mathbf{K}\) 上的一个有限维左向量空间。用 \(\mathfrak{N}(\mathbf{E})\) 表示 \(\mathbf{E}\) 上全体范数所成的集，它是 \(\mathcal{C}(\mathbf{E};\mathbf{R})\)（由 \(\mathbf{E}\) 到 \(\mathbf{R}\) 中的（在典范拓扑下）连续映射组成的空间）的一个子空间。
\end{enumerate}

\begin{enumerate}
\def\labelenumi{\alph{enumi})}
\item
  当给 \(\mathcal{C}(\mathrm{E};\mathbf{R})\) 赋以紧收敛拓扑 \(*\)（对此拓扑它是一个 Fréchet 空间）时，集合 \(\mathfrak{N}(\mathrm{E})\) 在 \(\mathcal{C}(\mathrm{E};\mathbf{R})\) 中是闭的，并且是局部紧的。
\item
  设 \(p_0\) 是 \(\mathfrak{N}(\mathrm{E})\) 的一个元素；证明：存在连续映射 \((\lambda, p) \mapsto \pi_{\lambda}(p)\)（\([0, 1] \times \mathfrak{N}(\mathrm{E})\) 到 \(\mathfrak{N}(\mathrm{E})\) 中的），使得 \(\pi_0(p) = p\) 且 \(\pi_1(p) = p_0\)（对每个 \(p \in \mathfrak{N}(\mathrm{E})\)）。
\end{enumerate}

\begin{enumerate}
\def\labelenumi{\arabic{enumi})}
\setcounter{enumi}{10}
\tightlist
\item
  在 I, 第 23 页, 习题 7 的假设下，证明：若 K 与 E 都完备，则 E 的每个闭子空间都有拓扑补（按同处 a) 的方式处理）。
\end{enumerate}

¶ 12) 设 K 是一个局部紧的赋值除环，其绝对值是非离散且超度量的。我们把 K 上左向量空间 E 上满足超度量不等式 (II, 第 2 页) 的范数称为超范数。

\begin{enumerate}
\def\labelenumi{\alph{enumi})}
\tightlist
\item
  设 E 是 K 上的一个有限维左向量空间，\(\alpha\) 是 E 上的一个超范数，H 是 E 中由方程 \(\langle x, a^{*} \rangle = 0\) 给出的一个超平面。证明：存在一点 \(x_{0} \in E\)，函数 \(x \mapsto |\langle x, a^{*} \rangle| / \alpha(x)\) 在 \(E \setminus \{0\}\) 上于该点达到其上确界；证明此时
\end{enumerate}

\[
\alpha (x) = \sup \left(\alpha \left(x - \frac {\langle x , a ^ {*} \rangle}{\langle x _ {0} , a ^ {*} \rangle} x _ {0}\right), \frac {| \langle x , a ^ {*} \rangle |}{| \langle x _ {0} , a ^ {*} \rangle |} \alpha (x _ {0})\right).
\]

由此推出：存在 E 的一个基 \((a_{i})\) 以及一个实数的族 \((r_i)\)（均 \(>0\)），使得对所有 \(x = \sum_{i} \xi_{i} a_{i}\)，有 \(\alpha(x) = \sup_{i} (r_{i} |\xi_{i}|)\)。我们称 \(\alpha\) 关于基 \((a_{i})\) 具有标准形式。

\begin{enumerate}
\def\labelenumi{\alph{enumi})}
\setcounter{enumi}{1}
\tightlist
\item
  设 \(\alpha^{*}\) 是 E 的对偶 E* 上由 \(\alpha\) 通过下式典范地确定的范数
\end{enumerate}

\[
\alpha^ {*} (x ^ {*}) = \sup _ {x \neq 0} | \langle x, x ^ {*} \rangle | / \alpha (x);
\]

它是一个超范数。证明：对所有 \(x_0 \neq 0\)（在 \(\mathbf{E}\) 中），存在 \(x_0^* \in \mathbf{E}^*\)，使得 \(\alpha(x_0) = |\langle x_0, x_0^* \rangle| / \alpha^*(x_0^*)\)。

\begin{enumerate}
\def\labelenumi{\alph{enumi})}
\setcounter{enumi}{2}
\tightlist
\item
  设 \(\alpha, \beta\) 是 E 上任意两个超范数。证明：存在 E 的一个基，使 \(\alpha\) 与 \(\beta\) 关于该基都具有标准形式（考察一点 \(x_0 \in \mathrm{E} \setminus \{0\}\)，\(\alpha / \beta\) 在该点达到其最大值；然后利用 \(b\)），并对 \(\dim \mathrm{E}\) 施行归纳法）。
\end{enumerate}

\(d)\) 把 \(\mathfrak{N}_0(\mathrm{E})\)（\(\mathrm{E}\) 上全体超范数所成的集）视为 \(\mathfrak{N}(\mathrm{E})\)（习题 10）的子空间。证明 \(\mathfrak{N}_0(\mathrm{E})\) 在 \(\mathfrak{N}(\mathrm{E})\) 中是闭的。设 \(\alpha_0\) 是 \(\mathfrak{N}_0(\mathrm{E})\) 的一个元素；对每个 \(\alpha \in \mathfrak{N}_0(\mathrm{E})\) 以及 \(0 \leqslant t \leqslant 1\)，用 \(\mathrm{P}_{\alpha}(t)\) 表示诸 \(\beta \in \mathfrak{N}_0(\mathrm{E})\)（满足 \(\beta(x) \leqslant \alpha_0(x)^{1 - t}\alpha(x)^t\)，对所有 \(x \in \mathrm{E}\)）组成的集。证明 \(\mathrm{P}_{\alpha}(t)\) 非空，且 \(\pi_t^\alpha = \sup \mathrm{P}_{\alpha}(t)\) 是一个超范数。此外，映射 \((t, \alpha) \mapsto \pi_t^\alpha\)（\([0, 1] \times \mathfrak{N}_0(E)\) 到 \(\mathfrak{N}_0(E)\) 中的）是连续的，且满足 \(\pi_0^\alpha = \alpha_0\) 与 \(\pi_1^\alpha = \alpha\)（利用 \(c\)）。

* e) 设 A 是 K 的绝对值的环，m 是其极大理想，使得 \(k = A/m\) 是一个含 q 个元素的有限域 (CA, VI, § 5, No.~1, prop. 2)。对 E 上每个超范数 \(\alpha\)，记 \(X_{\alpha}\) 为 \(\log \alpha(x)\)（\(x \in E \setminus \{0\}\)）的值集在到商群 \(\mathbf{R}/(\mathbf{Z}, \log q)\) 中的典范映射下的像；它是一个至多含 \(n = \dim E\) 个元素的有限集（利用 a)）；这些元素的个数记作 \(r(\alpha)\)，称为 \(\alpha\) 的秩。证明 r 是 \(\mathfrak{N}_{0}(E)\) 到 N 中的一个下半连续映射，且集 \(\mathfrak{N}_{0}'(E)\)（由诸 \(\alpha\)（\(r(\alpha) = n\)）所组成）在 \(\mathfrak{N}_{0}(E)\) 中是开的且处处稠密（利用 a) 与 c)）。

\begin{enumerate}
\def\labelenumi{\alph{enumi})}
\setcounter{enumi}{5}
\tightlist
\item
  设 \(r(\alpha) = n\)；设 \((a_i)\) 是 E 的一个基，\(\alpha\) 关于它具有标准形式；证明：存在 \(\alpha\) 在 \(\mathfrak{N}_0'(\mathrm{E})\) 中的一个邻域 V，使得每个 \(\beta \in V\) 都具有关于 \((a_i)(\text{使用 } b)\) 的标准形式；由此推出，存在邻域 \(W \subset V\)（\(\alpha\) 的），它同胚于 \(R^n\) 中的一个开集。
\end{enumerate}

\(g)\) 对 E 的每个基 \((a_i)\)，证明：超范数 \(\alpha\)（具有关于 \((a_i)\) 的标准形式）所成的集在 \(\mathfrak{N}_0^{\prime}(\mathrm{E})\) 中是闭的。由此推出：若 \(\alpha \in \mathfrak{N}_0^{\prime}(\mathrm{E})\) 具有关于 \((a_i)\) 的标准形式，则含 \(\alpha\) 的连通分支（\(\mathfrak{N}_0^{\prime}(\mathrm{E})\) 中的）的所有元素亦然。

¶ 13) * 我们保留习题 12 的一般假设与记号。

\begin{enumerate}
\def\labelenumi{\alph{enumi})}
\item
  设 \(\mathbf{L}\) 是 \(\mathbf{E}\) 的一个维数为 \(n = \dim \mathbf{E}\) 的自由 A-子模。对所有 \(x \in \mathbf{E} \setminus \{0\}\)，诸 \(a \in \mathbf{A}\)（满足 \(ax \in \mathbf{L}\)）组成的集是 \(\mathbf{K}\) 的一个形如 \(\mathfrak{m}^h\) 的分式理想（\(h\) 为正整数或负整数）；置 \(\alpha(x) = q^h\) 及 \(\alpha(0) = 0\)，证明 \(\alpha\) 是 \(\mathbf{E}\) 上的一个超范数。称它为与自由 A-模 \(\mathbf{L}\) 相伴的超范数。
\item
  反之，若 \(\alpha\) 是 E 上的一个超范数，则集 \(L_{\alpha}\)（由诸 \(x \in E\)（满足 \(\alpha(x) \leqslant 1\)）组成）是一个维数为 \(n\) 的自由 A-模。若 \([\alpha]\) 是与 \(L_{\alpha}\) 相伴的范数，则有 \(\alpha \leqslant [\alpha] \leqslant q\alpha\)，且 \([\alpha]\) 是与自由 A-模相伴且 \(\geqslant \alpha\) 的诸范数的下确界。我们有 \([q\alpha] = q \cdot [\alpha]\)，且 \(\alpha(x) = \inf(q^{-t}[q^t\alpha](x))\)（对所有 \(x \in E\)），其中 \(t\) 取遍区间 {[}0, 1{]}。此外，函数 \(t \mapsto [q^t\alpha](x)\) 在该区间上是左连续的。
\item
  沿用同样的记号，证明：对 \(0 \leqslant t \leqslant 1\)，诸 \([q^{t}\alpha]\) 中至多有 n 个互不相同的超范数。反之，设 L 是与维数为 n 的自由 A-模相伴的超范数所成的集，\((\alpha_{t})_{0 \leqslant t \leqslant 1}\) 是 L 的一个递增的超范数族，且满足 \(\alpha_{1} = q\alpha_{0}\)。证明：存在 E 的一个基，使所有 \(\alpha_{t}\) 都关于它具有标准形式（若 \(u \in A\) 是赋值为 1 的一个元素，\(L_{t}\) 是诸 \(x \in E\)（满足 \(\alpha_{t}(x) \leqslant 1\)）组成的自由 A-模，则考察 k 上的向量空间 \(L_{t}/uL_{0}\)）。进一步推出：若对所有 \(x \in E\)，\(t \mapsto \alpha_{t}(x)\) 在 \([0, 1]\) 上左连续，则存在唯一的超范数 \(\alpha\)，使得 \(\alpha_{t} = [q^{t}\alpha]\)（对所有 \(t \in [0, 1]\)）。
\end{enumerate}

\(d\) ) 线性群 \(\mathbf{GL}(\mathrm{E})\) 连续地作用于 \(\mathfrak{N}_0(\mathrm{E})\)；证明它是真作用。对所有 \(\alpha \in \mathfrak{N}_0(\mathrm{E})\)，稳定子 \(S_{\alpha}\)（\(\alpha\) 在 \(\mathbf{GL}(\mathrm{E})\) 中的）是诸 \([q^t\alpha]\)（\(0\leqslant t\leqslant 1\)）的稳定子的交；由此推出 \(S_{\alpha}\) 是 \(\mathbf{GL}(\mathrm{E})\) 的一个开紧子群，从而每个 \(\alpha \in \mathfrak{N}_0(\mathrm{E})\) 的轨道都是 \(\mathfrak{N}_0(\mathrm{E})\) 的一个闭离散子空间。

\begin{enumerate}
\def\labelenumi{\alph{enumi})}
\setcounter{enumi}{4}
\tightlist
\item
  对每个超范数 \(\alpha \in \mathfrak{N}_0(\mathrm{E})\)，考察向量 \(k\)-空间 \(\mathrm{L}_t / u\mathrm{L}_0\) 的维数所成的递减序列，其中 \(\mathrm{L}_t\) 是诸 \(x \in \mathrm{E}\)（满足 \([q^t\alpha](x) \leqslant 1\)）组成的 A-模，而 \(t\) 从 0 变到 1；我们称该序列为 \(\alpha\) 的不变量序列。为使 \(\alpha\) 与 \(\beta\) 属于 \(\mathfrak{N}_0(\mathrm{E})\) 中同一轨道，必须且只需 \(\mathrm{X}_{\alpha} = \mathrm{X}_{\beta}\)（习题 12, e)）且 \(\alpha\) 与 \(\beta\) 的不变量序列相同（利用习题 12, b)）。f) 由 e) 推出：轨道空间 \(\mathfrak{N}_0(\mathrm{E}) / \mathbf{GL}(\mathrm{E})\) 同构于轨道空间 \(\mathbf{T}^n / \mathfrak{S}_n\)，其中对称群在 \(\mathbf{T}^n\) 上通过 \((z_1, ..., z_n) \mapsto (z_{\sigma(1)}, ..., z_{\sigma(n)})\) 右作用。*
\end{enumerate}

\begin{enumerate}
\def\labelenumi{\arabic{enumi})}
\setcounter{enumi}{13}
\item
  把第 2 小节与第 3 小节的结果推广到离散除环 K 上的这样的拓扑向量空间 E：E 中存在由均衡邻域组成的 0 的基本邻域系（即由满足 K.V = V 的邻域 V 组成的基本邻域系）。
\item
  设 \(\mathbf{E}\) 是一个维数 \(n\) 有限的赋范空间（在 \(\mathbf{R}\) 或 \(\mathbf{C}\) 上）。给对偶 \(\mathbf{E}^*\) 赋以由 \(\| x^* \| = \sup_{\| x \| \leqslant 1} |\langle x, x^* \rangle|\) 定义的范数 (GT, X, § 3.2)。证明：存在一个基 \((e_i)\)
\end{enumerate}

使 E 满足：若 \((e_{i}^{*})\) 是对偶基，则对所有 i 有 \(\|e_{i}\| = \|e_{i}^{*}\| = 1\)。（设 \((a_{i})\) 是由范数为 1 的向量组成的 E 的一个基；考察行列式 \(\det(\xi_{ij})\)（对每个由 n 个范数为 1 的向量 \(x_{i} = \sum_{j} \xi_{ij} a_{j}\) 组成的向量组），并考察使该行列式的绝对值达到最大的这样一个向量组。）

\exercisegroup

\begin{enumerate}
\def\labelenumi{\arabic{enumi})}
\tightlist
\item
  \begin{enumerate}
  \def\labelenumii{\alph{enumii})}
  \tightlist
  \item
    证明：若非离散赋值除环 K 上的 Hausdorff 拓扑向量空间 E 使得 0 的每个邻域都包含一个并非单点 0 的向量子空间，则 E 的拓扑不能由范数定义。特别地，由 K 上的 Hausdorff 拓扑向量空间组成的无穷序列 \((\mathrm{E}_{n})\)（其中任何一个都不由单点 0 组成）的乘积具有一个不能由范数定义的拓扑。
  \end{enumerate}
\end{enumerate}

\begin{enumerate}
\def\labelenumi{\alph{enumi})}
\setcounter{enumi}{1}
\tightlist
\item
  考察乘积向量空间 \(E = K_{s}^{N}\)；对所有 \(x = (\xi_{n}) \in E\)，置
\end{enumerate}

\[
| x | = \sum_ {n = 0} ^ {\infty} 2 ^ {- n} \left| \xi_ {n} \right| / (1 + | \xi_ {n} |).
\]

证明：E 的拓扑由距离 \(d(x, y) = |x - y|\) 定义；\(|\lambda x| \leqslant |x|\)（若 \(|\lambda| \leqslant 1\)），\(|\lambda x| \leqslant |\lambda| \cdot |x|\)（若 \(|\lambda| \geqslant 1\)）；并且对所有 \(x_0 \in \mathrm{E}\)，\(|\lambda x_0|\) 随 \(|\lambda|\) 趋于 0。

\begin{enumerate}
\def\labelenumi{\arabic{enumi})}
\setcounter{enumi}{1}
\tightlist
\item
  设 E 与 F 是非离散赋值除环上的两个完备可度量化向量空间，\(T_{0}\) 是 F 的拓扑。设 T 是 F 上的一个 Hausdorff 拓扑，比 \(T_{0}\) 更粗。证明：若 E 到 F 中的线性映射 u 关于 F 上的拓扑 T 连续，则它关于 F 上的拓扑 \(T_{0}\) 仍连续（利用 I, 第 19 页 的 cor. 5）。
\end{enumerate}

由此推出：若 \(T_{1}\) 与 \(T_{2}\) 是非离散赋值除环上的向量空间 E 上两个不同的拓扑，它们都与 E 的向量空间结构相容，且对其中的每一个，E 都是可度量化和完备的，则 E 上不存在比 \(T_{1}\) 与 \(T_{2}\) 都更粗的 Hausdorff 拓扑。在无穷维向量空间 E 上给出这样两个拓扑的例子（注意：存在 E 到自身的双射，使得该双射及其逆映射关于 E 上的一个赋范空间拓扑都不连续）。

\begin{enumerate}
\def\labelenumi{\arabic{enumi})}
\setcounter{enumi}{2}
\item
  设 E 与 F 是非离散赋值除环上的两个 Hausdorff 拓扑向量空间，并假设 E 是可度量化和完备的。设 u 是 E 到 F 中的一个连续线性单射，G 是 \(u(\mathrm{E})\) 的一个向量子空间；假设 G 上存在一个拓扑 T，它比由 F 的拓扑诱导的拓扑更细，与 G 的向量空间结构相容，且对此拓扑 G 是可度量化和完备的。证明：u 的逆映射限制在 G 上关于 T 连续（利用 I, 第 19 页, cor. 5）。
\item
  设 E、F 是非离散赋值除环上的两个完备可度量化向量空间，u 是 E 到 F 中的一个连续线性映射。证明：若 F 中存在 \(u(\mathrm{E})\) 的一个闭的补子空间，则 \(u(\mathrm{E})\) 在 F 中是闭的（利用 I, 第 19 页, cor. 5）。
\item
  设 E 与 F 是非离散赋值除环上的两个完备可度量化向量空间，u 是 E 到 F 中的一个线性映射。用 N 表示 u 在 F 中关于 E 中 0 的邻域滤子的接触点所成的集；证明 N 是 F 的一个闭向量子空间，并且为使 u 连续，必须且只需 N 是单点 0（利用 I, 第 19 页, cor. 5）。证明：在 F 的闭向量子空间 M 之中，N 是满足下列条件的最小者：若 \(\phi\) 表示 F 到 F/M 上的典范同态，则 \(\phi \circ u\) 是 E 到 F/M 中的一个连续映射。
\item
  设 E 是非离散赋值除环 K 上的一个完备可度量化向量空间。
\end{enumerate}

\begin{enumerate}
\def\labelenumi{\alph{enumi})}
\tightlist
\item
  设 p 是 E 的一个到区间 \([0, +\infty)\)（\(\overline{R}\) 中的）内的下半连续映射，满足 \(p(\lambda x) = |\lambda| \cdot p(x)\)（对 K 中每个 \(\lambda \neq 0\) 及每个 \(x \in E\)）；\(p(0) = 0\)；并且 \(p(x + y) \leqslant p(x) + p(y)\)
\end{enumerate}

对 E 中任何 \(x, y\) 成立。证明：若 \(p\) 在 E 中是有限的，则 \(p\) 连续（考察由诸 \(x \in \mathbf{E}\)（满足 \(p(x) \leqslant 1\)）组成的闭集 B，并利用 Baire 定理）。

\begin{enumerate}
\def\labelenumi{\alph{enumi})}
\setcounter{enumi}{1}
\tightlist
\item
  设 \((p_n)\) 是 \(E\) 到 \([0, +\infty]\) 中的映射序列，满足 \(a\) ) 中的条件。证明：若诸 \(p_n\) 中没有一个在 \(E\) 中是有限的，则存在一点 \(x \in E\)，使得 \(p_n(x) = +\infty\)（对所有 \(n\)）（方法同上）。
\end{enumerate}

\begin{enumerate}
\def\labelenumi{\arabic{enumi})}
\setcounter{enumi}{6}
\tightlist
\item
  设 E 是非离散赋值除环 K 上的一个完备可度量化向量空间。我们称 E 的向量子空间 M 是仿完备的，如果 M 上存在一个完备可度量化的向量空间结构，使得 M 到 E 中的典范单射连续。a) 设 M、N 是 E 的两个仿完备子空间，且 \(M + N\) 与 \(M \cap N\) 在 E 中都是闭的。证明 M 与 N 在 E 中是闭的。（取关于 \(M \cap N\) 的商可将问题化为 \(M \cap N = \{0\}\) 的情形，这时可以考察映射 \((x, y) \mapsto x + y\)（\(M \times N\) 到 E 中的）。）
\end{enumerate}

\begin{enumerate}
\def\labelenumi{\alph{enumi})}
\setcounter{enumi}{1}
\tightlist
\item
  证明：若 E 是仿完备子空间的一个递增序列 \((\mathbf{M}_{j})_{j\geq0}\) 的并，则存在一个指标 j，使得 \(M_{j}=E\)。（利用 Baire 定理 (GT, IX, § 5.3, th. 1) 及 I, 第 17 页 的 th. 1。）
\end{enumerate}

\begin{enumerate}
\def\labelenumi{\arabic{enumi})}
\setcounter{enumi}{7}
\tightlist
\item
  设 E 是非离散赋值除环 K 上的一个 Banach 空间。我们称 E 的向量子空间 M 是强仿完备的，如果 M 上存在一个范数 \(\|x\|_{M}\)，使得 M 关于它是一个 Banach 空间，且 M 到 E 中的典范单射连续。
\end{enumerate}

\begin{enumerate}
\def\labelenumi{\alph{enumi})}
\item
  证明：若 \(\mathbf{M}\) 与 \(\mathbf{N}\) 是 \(\mathbf{E}\) 的两个强仿完备子空间，则 \(\mathbf{M} + \mathbf{N}\) 与 \(\mathbf{M} \cap \mathbf{N}\) 也是强仿完备子空间。（在 \(\mathbf{M} + \mathbf{N}\) 上考虑范数 \(\|x\|_{\mathbf{M} + \mathbf{N}} = \inf (\|u\|_{\mathbf{M}} + \|v\|_{\mathbf{N}})\)，其中下确界取遍所有对 \((u, v)\)（满足 \(x = u + v\)、\(u \in \mathbf{M}\)、\(v \in \mathbf{N}\)）。）
\item
  设 M、N 是 E 的两个强仿完备子空间，且 N 与 \(M + N\) 都是闭的。证明 \(\overline{M} = M + (\overline{M} \cap \overline{N})\) 及 \(\overline{M} \cap \overline{N} = \overline{M} \cap N\)（利用 exerc. 7, a)）。
\end{enumerate}

\begin{enumerate}
\def\labelenumi{\arabic{enumi})}
\setcounter{enumi}{8}
\tightlist
\item
  \begin{enumerate}
  \def\labelenumii{\alph{enumii})}
  \tightlist
  \item
    设 \(a, b\) 是赋范空间 \(\mathbf{E}\)（在域 \(\mathbf{R}\) 上）中的两点。用 \(\delta(\mathbf{A})\) 表示有界集 \(\mathbf{A}\)（在 \(\mathbf{E}\) 中）的直径（使用 \(\mathbf{E}\) 上由范数导出的距离），并归纳地定义有界集的序列 \((\mathbf{B}_n)_{n \geqslant 1}\)（在 \(\mathbf{E}\) 中），使其满足下列条件：\(\mathbf{B}_1\) 是诸 \(x \in \mathbf{E}\)（满足 \(\|x - a\| = \|x - b\| = \frac{1}{2}\|a - b\|\)）组成的集；当 \(n > 1\) 时，\(\mathbf{B}_n\) 是诸 \(x \in \mathbf{B}_{n-1}\)（满足 \(\|x - y\| \leqslant \frac{1}{2}\delta(\mathbf{B}_{n-1})\)，对所有 \(y \in \mathbf{B}_{n-1}\)）组成的集。证明诸 \(\mathbf{B}_n\) 的交恰好是单点 \(\frac{1}{2}(a + b)\)（注意 \(\delta(\mathbf{B}_n) \leqslant \frac{1}{2}\delta(\mathbf{B}_{n-1})\)）。
  \end{enumerate}
\end{enumerate}

\begin{enumerate}
\def\labelenumi{\alph{enumi})}
\setcounter{enumi}{1}
\tightlist
\item
  由 \(a)\) 推出：若 \(u\) 是实 Banach 空间 \(E\) 到实 Banach 空间 \(F\) 上的一个等距，则 \(u\) 是 \(E\) 到 \(F\) 上的一个仿射线性映射。
\end{enumerate}

\chapter{凸集与局部凸空间}

在本章的 §§ 2 至 7 中，我们只考虑实数域 R 上的向量空间与仿射空间；当我们谈到一个向量空间或一个仿射空间而不明确给出其标量除环时，应理解该除环就是域 R。关于 C 上的向量空间，见 § 8。

\section{半范数}

在本节中，K 表示一个非离散赋值除环。

\subsection{半范数的定义}

定义 1. --- 设 E 是 K 上的一个左向量空间。E 到 \(R_{+} = \left[0, +\infty\right]\) 中的映射 p 称为 E 上的半范数，如果它满足下列公理：

\((\mathrm{SN}_{\mathrm{I}})\) 若 \(x \in \mathrm{E}\) 且 \(\lambda \in \mathrm{K}\)，则 \(p(\lambda x) = |\lambda| p(x)\)。

\((\mathrm{SN}_{\mathrm{II}})\) 若 \(x \in \mathrm{E}\) 且 \(y \in \mathrm{E}\)，则 \(p(x + y) \leqslant p(x) + p(y)\)。

由于 \(p(x) \leqslant p(y) + p(x - y)\) 及 \(p(y) \leqslant p(x) + p(y - x)\)，由 \(p(y - x) = p(x - y)\)，我们推出

\[
\left| p (x) - p (y) \right| \leqslant p (x - y).\tag{1}
\]

例. --- 1) E 上的范数是这样一个半范数 p：关系 \(p(x) = 0\) 蕴含 x = 0（I, 第 3 页）。

\begin{enumerate}
\def\labelenumi{\arabic{enumi})}
\setcounter{enumi}{1}
\item
  对每个线性型 \(f\)（在 \(\mathbf{E}\) 上），函数 \(x \mapsto |f(x)|\) 是 \(\mathbf{E}\) 上的半范数。
\item
  若 \(p_i (1 \leqslant i \leqslant n)\) 是 \(E\) 上有限个半范数，则显然 \(p'(x) = \sup_{1 \leqslant i \leqslant n} p_i(x)\)。
\end{enumerate}

与 \(p''(x) = \sum_{i=1}^{n} \alpha_i p_i(x)\)（其中诸 \(\alpha_i\) \(\geqslant 0\)）都是 E 上的半范数。

\(p\) 是 \(\mathbf{E}\) 到 \(\mathbf{R}_{+}\) 中的映射，若它满足 \((\mathrm{SN}_{\mathrm{I}})\) 及下列公理，则称为超半范数：

\((\mathrm{SN}_{\mathrm{II}}^{\prime})\) 若 \(x \in \mathrm{E}\) 且 \(y \in \mathrm{E}\)，则 \(p(x + y) \leqslant \sup(p(x), p(y))\)。

显然，超半范数是半范数。

说 K 上的绝对值是超度量的 (CA, VI, § 6.2)，就是指它是左向量空间 \(K_{s}\) 上一个不恒为零的超半范数。

命题 1. --- 设 E 是 K 上的一个左拓扑向量空间，p 是 E 上的一个半范数。下列条件等价：

\begin{enumerate}
\def\labelenumi{\alph{enumi})}
\item
  \(p\) 在 E 中连续。
\item
  \(p\) 在点 0 连续。
\item
  \(p\) 一致连续。
\item
  对每个实数 \(\alpha > 0\)，集 \(\mathbf{W}(p, \alpha)\)（诸 \(x \in \mathbf{E}\)（满足 \(p(x) < \alpha\)）所组成的）在 \(\mathbf{E}\) 中是开的。
\item
  存在实数 \(\alpha > 0\)，使得 \(\mathrm{W}(p, \alpha)\) 是 E 中 0 的一个邻域。
\item
  对每个实数 \(\alpha > 0\)，集 \(\mathrm{V}(p, \alpha)\)（诸 \(x \in \mathrm{E}\)（满足 \(p(x) \leqslant \alpha\)）所组成的）是 \(\mathrm{E}\) 中 0 的一个邻域。
\end{enumerate}

事实上，蕴含关系 \(c) \Rightarrow a) \Rightarrow b) \Rightarrow d) \Rightarrow e) \Rightarrow f) \Rightarrow c)\) 由 \((\mathrm{SN}_{\mathrm{I}})\) 及不等式 (1) 立即得出。

推论. --- 若 p 是 E 上的一个连续半范数，q 是满足 \(q \leqslant p\) 的半范数，则 q 在 E 中连续。

当 \(p\) 是 \(\mathbf{E}\) 上的超半范数时，集 \(\mathrm{W}(p, \alpha)\) 与 \(\mathrm{V}(p, \alpha)\) 都是既开又闭的。因为，我们已看到 \(\mathrm{W}(p, \alpha)\) 是开的；另一方面，若 \(z\) 是 \(\mathrm{W}(p, \alpha)\) 的一个接触点，则存在 \(y \in \mathrm{W}(p, \alpha)\) 使得 \(p(y - z) < \alpha\)，而由 \((\mathrm{SN}_{\mathrm{II}}')\) 有 \(p(z) < \alpha\)，故 \(\mathrm{W}(p, \alpha)\) 是闭的。又，\(\mathrm{V}(p, \alpha)\) 是闭的（因为 \(p\) 连续）；进而，若 \(p(x) \leqslant \alpha\) 且 \(p(y) \leqslant \alpha\)，则 \(p(x + y) \leqslant \alpha\)（由 \((\mathrm{SN}_{\mathrm{II}}')\)），这就表明 \(\mathrm{V}(p, \alpha)\) 是开的。

\subsection{由半范数定义的拓扑}

设 \(p\) 是向量空间 \(E\)（在 \(K\) 上）上的一个半范数；对每个 \(\alpha > 0\)，用 \(V(p, \alpha)\) 表示由那些 \(x\)（属于 \(E\) 且满足 \(p(x) \leqslant \alpha\)）组成的子集。显然，若 \(x \in V(p, \alpha)\) 且 \(\lambda \in K\) 满足 \(|\lambda| \leqslant 1\)，则 \(\lambda x \in V(p, \alpha)\)，换言之，\(V(p, \alpha)\) 是均衡的。又，对每个 \(x_0 \in E\)，存在非零标量 \(\mu \in K\) 使得 \(|\mu| \geqslant p(x_0) \alpha^{-1}\)，因此 \(\mu^{-1}x_0 \in V(p, \alpha)\)，也就是说 \(V(p, \alpha)\) 是吸收的。最后，由 \((\mathrm{SN}_{\mathrm{II}})\) 有 \(V(p, \alpha/2) + V(p, \alpha/2) \subset V(p, \alpha)\)，并由 \((\mathrm{SN}_{\mathrm{I}})\)，对每个非零标量 \(\lambda\)（在 \(K\) 中）有 \(\lambda V(p, \alpha) = V(p, |\lambda| \alpha)\)。由这些备注并根据 I, 第 7 页 的 prop. 4，我们断定：当 \(\alpha\) 在 \(> 0\) 的数集中变化（或只在趋于 0 的一个严格正数序列中变化）时，诸集 \(V(p, \alpha)\) 构成 \(E\) 的一个与向量空间结构相容的拓扑的 0 的基本邻域系；我们称这拓扑由半范数 \(p\) 定义。赋有这种拓扑的向量空间 \(E\) 称为半赋范空间。注意，若 \(W(p, \alpha)\) 表示由那些 \(x\)（属于 \(E\) 且满足 \(p(x) < \alpha\)）组成的子集，则诸 \(W(p, \alpha)\)（其中 \(\alpha > 0\)，或 \(\alpha\)）

在一个趋于零的严格正数序列中变化）构成由 \(p\) 定义的拓扑的 0 的一个基本邻域系。

若 \(\Gamma\) 是 \(\mathbf{E}\) 上半范数的一个集合，则由诸半范数 \(p \in \Gamma\) 定义的拓扑的上确界与向量空间结构相容（I, 第 10 页, cor. 4）。这个拓扑的 0 的一个基本邻域系由诸有限交 \(\bigcap_{i} \mathrm{V}(p_i, \alpha_i)\) 给出，其中 \(p_i \in \Gamma\) 且 \(\alpha_i > 0\)。称这拓扑是由半范数集

\(\Gamma\) 定义的。它是 \(\mathbf{E}\) 上在所有平移下不变且使诸半范数 \(p \in \Gamma\) 连续的拓扑中最粗的一个。

设 E 是 K 上的一个拓扑向量空间：E 上的半范数组 \(\Gamma\) 称为基本半范数系，如果 E 上的拓扑与由 \(\Gamma\) 定义的拓扑相同。

设 E 是 K 上的向量空间，赋有由半范数集 \(\Gamma\) 定义的拓扑。对每个半范数 p，有 \(p(x - z) \leqslant p(x - y) + p(y - z)\)，这表明函数 \((x, y) \mapsto p(x - y)\) 是 E 上的一个伪距离 (GT, IX, §1.1)：由定义可得，当 p 在 \(\Gamma\) 中变化时，这些伪距离组成的集定义了拓扑向量空间 E 的一致结构。

注记. --- 1) E 上由有限个半范数 \(p_i\)（\(1 \leqslant i \leqslant n\)）定义的拓扑可由单个半范数 \(p = \sup_{1 \leqslant i \leqslant n} p_i\) 定义。但由

无穷多个半范数定义的拓扑一般不能由单个半范数定义（III, 第 37 页, exerc. 2）。

\begin{enumerate}
\def\labelenumi{\arabic{enumi})}
\setcounter{enumi}{1}
\item
  设 \((\mathcal{T}_{\iota})_{\iota\in I}\) 是 K 上向量空间 E 上的一族拓扑，其中每个拓扑都由一族半范数 \(\Gamma_{\iota}\) 定义。则由半范数集 \(\Gamma = \bigcup \Gamma_{\iota}\) 定义的拓扑是诸拓扑 \(T_{\iota}\) 的上确界。
\item
  若 \(\Gamma_0\) 是一个半范数集，它由 E 上两个半范数 \(p, q\) 之间如下定义的递增序关系所定向：「存在 \(\lambda > 0\) 使得 \(p \leqslant \lambda q\)」，则对由 \(\Gamma_0\) 定义的拓扑，取诸集 \(\mathbf{V}(p, \alpha)\)（其中 \(p \in \Gamma_0\)，\(\alpha > 0\)）即得 0 的一个基本邻域系。若 \(\Gamma\) 是 E 上任意一个半范数集，则与 \(\Gamma\) 定义相同拓扑的一个滤过化半范数集，就是集 \(\Gamma_0\)（由属于 \(\Gamma\) 的半范数的一切有限族的上包络所组成）。
\item
  即使 \(\mathbf{K} = \mathbf{R}\)，\(\mathbf{K}\) 上拓扑向量空间的拓扑也不总能由半范数集定义（cf.~II, 第 24 页, corollary）。
\end{enumerate}

例. --- 设 \(\mathcal{C}^{\infty}(\mathbf{R})\) 是 R 上在 R 中无穷次可微的实值函数组成的 R 上的向量空间。对每个函数以及每对整数 \(n \geqslant 0\)、\(m \geqslant 1\)，置

\[
p _ {n, m} (f) = \sup _ {- m \leqslant t \leqslant m} | f ^ {(n)} (t) |\tag{2}
\]

其中 \(f^{(0)} = f\)。显然诸 \(p_{n,m}\) 是 \(\mathcal{C}^{\infty}(\mathbf{R})\) 上的半范数。为使诸函数 \(f_{\alpha}\) 趋于 0（按指标集上的一个滤子 \(\mathfrak{F}\)），在 \(\mathcal{C}^{\infty}(\mathbf{R})\) 中关于拓扑 \(\mathcal{T}\)（由半范数 \(p_{n,m}\) 定义的），必须且只需：对所有整数 \(n \geqslant 0\)，诸函数 \(f_{\alpha}^{(n)}\)（按 \(\mathfrak{F}\)）在 \(\mathbf{R}\) 的每个紧子集上一致趋于 0。我们称 \(\mathcal{T}\) 是函数 \(f \in \mathcal{C}^{\infty}(\mathbf{R})\) 及其所有导数的紧收敛拓扑（cf.~III, 第 9 页）。

命题 2. --- 在向量空间 E 上，设 T 是由半范数集 Γ 定义的拓扑。

\begin{enumerate}
\def\labelenumi{(\roman{enumi})}
\item
  \(\{0\}\) 在 \(\mathbf{E}\) 中关于 \(\mathcal{T}\) 的闭包，是由诸 \(x \in \mathbf{E}\)（满足 \(p(x) = 0\)，对每个半范数 \(p \in \Gamma\)）组成的子集。
\item
  若 \(\mathcal{T}\) 是 Hausdorff 的且 \(\Gamma\) 可数，则 \(\mathcal{T}\) 可度量化。
\end{enumerate}

本命题由定义及 GT, IX, § 2.4, cor. 1 立即得出。

注意，即使 \(\mathcal{T}\) 可度量化，\(\mathcal{T}\) 也可能不能由单个范数来定义；上面给出的例子正是这种情形（cf.~IV, 第 18 页, 例 4）。

设 E 是 K 上的向量空间，赋有由半范数集 \(\Gamma\) 定义的拓扑。设 \(\hat{E}\) 是 E 的 Hausdorff 完备化（I, 第 6 页），\(\hat{\Gamma}\) 是映射 \(\hat{p}\)（\(\hat{E}\) 到 \(R_{+}\) 中的）组成的集，其中 p 在 \(\Gamma\) 中变化（GT, II, § 3.7, prop. 15）。由不等式延拓原理，诸函数 \(\hat{p} \in \hat{\Gamma}\) 是 \(\hat{E}\) 上的半范数，且诸函数 \(\hat{p}(x - y)\) 组成一个定义 \(\hat{E}\) 的一致结构的伪距离集（GT, IX, § 1.3, prop. 1）。因此我们看到，\(\hat{\Gamma}\) 是定义 \(\hat{E}\) 的拓扑的一个基本半范数集。

\subsection{商空间与乘积空间中的半范数}

设 E 是 K 上的一个拓扑向量空间，其拓扑由半范数集 \(\Gamma\) 定义。显然，\(\Gamma\) 中诸半范数在 E 的向量子空间 M 上的限制，定义了由 E 的拓扑在 M 上诱导的拓扑。

设 \(\phi\) 是 E 到向量商空间 E/M 上的典范映射。我们证明，对 E 上的每个半范数 \(p\)，函数

\[
\dot {p} (z) = \inf _ {\phi (x) = z} p (x)\tag{3}
\]

是 E/M 上的半范数。事实上，显然 \(\dot{p}\) 满足条件 \((\mathrm{SN}_{\mathrm{I}})\)；另一方面，若 \(z'\)、\(z''\) 是 E/M 的两个向量，我们有：

\[
\begin{array}{r l} \inf _ {\phi (x) = z ^ {\prime} + z ^ {\prime \prime}} p (x) & \leqslant \inf _ {\phi (x ^ {\prime}) = z ^ {\prime}, \phi (x ^ {\prime \prime}) = z ^ {\prime \prime}} p (x ^ {\prime} + x ^ {\prime \prime}) \\ & \leqslant \inf _ {\phi (x ^ {\prime}) = z ^ {\prime}, \phi (x ^ {\prime \prime}) = z ^ {\prime \prime}} \left(p (x ^ {\prime}) + p (x ^ {\prime \prime})\right) \\ & = \inf _ {\phi (x ^ {\prime}) = z ^ {\prime}} p (x ^ {\prime}) + \inf _ {\phi (x ^ {\prime \prime}) = z ^ {\prime \prime}} p (x ^ {\prime \prime}) \end{array}
\]

这表明 \(\dot{p}\) 满足 \((\mathrm{SN}_{\mathrm{II}})\)。我们称 \(\dot{p}\) 是 p 关于 M 的商半范数。

同样的推理证明：若 p 是超半范数，则 \(\dot{p}\) 也是超半范数。

这样，我们有（采用第 2 目的记号）

\[
\phi (\mathrm{W} (p, \alpha)) = \mathrm{W} (\dot {p}, \alpha).\tag{4}
\]

对每个 \(\alpha > 0\) 成立。事实上，说 \(\dot{p}(z) < \alpha\)，意味着存在 \(x \in \mathbf{E}\) 使得 \(\phi(x) = z\) 且 \(p(x) < \alpha\)，由此即得关系式 (4)。

由此我们推出：若半范数集 \(\Gamma\) 是有向的（II, 第 3 页, 注记 3），则 E/M 上的商拓扑由诸半范数 \(\dot{p}\)（其中 p 在 \(\Gamma\) 中变化）所成的集定义。

若 N 是 \(\{0\}\) 在 E 中的闭包，则 E/N 的拓扑由诸商半范数 \(\dot{p}\)（其中 p 在 \(\Gamma\) 中变化）定义（即使 \(\Gamma\) 不是滤过的）：这里，\(\dot{p}(\dot{x}) = p(x)\)（对属于类 \(\dot{x} \mod N\) 的每个 x）。注意，E/N 不是别的，正是与 E 相伴的 Hausdorff 空间（I, 第 4 页）。

设 \(\mathbf{E}\) 是 \(\mathbf{K}\) 上的向量空间，\((\mathrm{E}_{\iota})_{\iota \in \mathbf{I}}\) 是 \(\mathbf{K}\) 上的一族向量空间，其中 \(\mathrm{E}_{\iota}\) 赋有拓扑 \(\mathcal{T}_{\iota}\)（由半范数集 \(\Gamma_{\iota}\) 定义的）。对每个 \(\iota \in \mathbf{I}\)，设 \(f_{\iota}\) 是 \(\mathbf{E}\) 到 \(\mathrm{E}_{\iota}\) 中的一个线性映射；显然，当 \(p_{\iota}\) 在集 \(\Gamma_{\iota}\) 中变化时，诸 \(p_{\iota} \circ f_{\iota}\) 构成一个半范数集 \(\Gamma_{\iota}'\)（\(\mathbf{E}\) 上的）。于是，拓扑 \(\mathcal{T}\)（\(\mathbf{E}\) 上的，定义为使所有映射 \(f_{\iota}\) 都连续的拓扑中最粗者；I, 第 9 页）由半范数集 \(\Gamma' = \bigcup_{\iota \in \mathbf{I}} \Gamma_{\iota}'\) 定义，这由 \(\mathcal{T}\) 的 0 的邻域的定义得出（GT, I, § 2.3, prop. 4）。

若诸 \(p_{\iota}\) 是超半范数，则诸 \(p_{\iota} \circ f_{\iota}\) 也是超半范数。

设 E 是 K 上的向量空间，赋有由半范数族 \((p_{\iota})_{\iota \in I}\) 定义的拓扑 T；对每个 \(\iota \in I\)，设 \(T_{\iota}\) 是由单个半范数 \(p_{\iota}\) 定义的拓扑，并用 \(E_{\iota}\) 表示由 E 赋有拓扑 \(T_{\iota}\) 而得的空间。则拓扑 T 是对角映射 \(\Delta : E \to \prod_{\iota \in I} E_{\iota}\) 下 \(\prod_{\iota \in I} \mathrm{E}_{\iota}\) 上乘积拓扑的原像（I, 第 9 页, prop. 7）。对每个 \(\iota \in I\)，用 \(\mathbf {N} _ {\iota}\) 表示 \(\{0 \}\) 在 \(\mathrm{E} _ {\iota}\) 中的闭包，并用 \(F_{\iota} = E_{\iota}/N_{\iota}\) 表示由范数 \(\hat{p}_{\iota}\)（对应于 \(p_{\iota}\) 的）定义的赋范空间（II, 第 4 页, 公式 (3)）；若 \(\phi_{\iota}: E_{\iota} \to F_{\iota}\) 是典范映射，\(\phi : (x_{\iota}) \mapsto (\phi_{\iota}(x_{\iota}))\) 是乘积映射，我们知道 \(\prod_{i \in I} E_{i}\) 上的乘积拓扑是 \(\phi\) 下 \(\prod_{\iota \in I}F_{\iota}\) 上乘积拓扑的原像（GT, II, § 3.9, prop. 18）。因此，拓扑 \(\mathcal{T}\) 是复合映射 \(\phi \circ \Delta\) 下 \(\prod_{\mathfrak{t}\in \mathbf{I}}\mathrm{F}_{\mathfrak{t}}\) 上乘积拓扑的原像。特别地，若 \(\mathcal{T}\) 是 Hausdorff 的，则由 II, 第 3 页, prop. 2 得 \(\Delta^{-1}(\prod_{\iota \in I} \mathbf{N}_{\iota}) = \{0 \}\)，因而映射 \(\phi \circ \Delta\) 是单射，因而：

命题 3. --- K 上每个其拓扑由半范数集定义的 Hausdorff 拓扑向量空间 E，都同构于一个 Banach 空间乘积的子空间。

进而，若 E 的拓扑由可数的半范数集定义，则 E 可度量化（I, 第 16 页）。

\subsection{多线性映射关于由半范数定义的拓扑的等度连续性判别法}

命题 4. --- 设 \(E_{i}\)（\(1 \leqslant i \leqslant n\)）与 F 是 K 上的拓扑向量空间；我们假设，对每个 i，\(E_{i}\) 的拓扑由一个有向半范数集 \(\Gamma_{i}\) 定义，而 F 的拓扑由半范数集 \(\Gamma\) 定义。那么，\(\prod_{i=1}^{n} \mathrm{E}_{i}\) 到 F 中的多线性映射所成的一个集 H 是等度连续的，当且仅当：对每个半范数 \(q \in \Gamma\) 及每个指标 \(i\)，存在一个半范数 \(p_{i} \in \Gamma_{i}\) 及一个数 \(a > 0\)，使得对每个映射 \(u \in \mathrm{H}\) 及每点 \((x_{i}) \in \prod_{i=1}^{n} \mathrm{E}_{i}\)，

\[
q \left(u \left(x _ {1}, x _ {2}, \dots , x _ {n}\right)\right) \leqslant a. p _ {1} \left(x _ {1}\right) p _ {2} \left(x _ {2}\right) \dots p _ {n} \left(x _ {n}\right).\tag{5}
\]

条件是充分的，因为它蕴含 H 在 \((0, 0, \ldots, 0)\) 处等度连续，从而处处等度连续（I, 第 9 页, prop. 6）。

我们证明条件是必要的。由假设，对每个半范数 \(q \in \Gamma\) 及每个数 \(\beta > 0\)，都有 \(q(u(x_1, x_2, ..., x_n)) \leqslant \beta\)（对每个映射 \(u \in \mathbf{H}\)），只要 \(p_i(x_i) \leqslant \alpha_i\) 对每个指标 \(i\)（\(1 \leqslant i \leqslant n\)）成立，而数 \(\alpha_i > 0\) 与半范数 \(p_i \in \Gamma_i\) 是适当选取的。由于 \(\mathbf{K}\) 非离散，我们还可假设对每个 \(i\)，\(\alpha_i = |\lambda_i| < 1\)，其中 \(\lambda_i \in \mathbf{K}\)。现在设 \((x_1, x_2, ..., x_n)\) 是 \(\prod_{i=1}^{n} E_i\) 中任意一点，并对每个指标 \(i\)，设 \(m_i \in \mathbf{Z}\) 是满足 \(p_i(x_i) \leqslant |\lambda_i|^{m_i + 1}\) 的整数；这可写成 \(p_i(\lambda_i^{-m_i} x_i) \leqslant |\lambda_i| (1 \leqslant i \leqslant n)\)，因此由假设，我们有

\[
q \left(u \left(x _ {1}, x _ {2}, \dots , x _ {n}\right)\right) \leqslant \beta \left| \lambda_ {1} \right| ^ {m _ {1}} \left| \lambda_ {2} \right| ^ {m _ {2}} \dots \left| \lambda_ {n} \right| ^ {m _ {n}}.\tag{6}
\]

首先假设某个 \(p_{i}(x_{i})\) 为零，那么可取 \(m_{i} \in N\) 任意大，因此 \(q(u(x_{1}, x_{2}, ..., x_{n})) = 0\)。反之，若所有 \(p_{i}(x_{i})\) 都 \(\neq 0\)，则对每个 i 取整数 \(m_{i}\) 使得 \(|\lambda_{i}|^{m_{i}+2} < p_{i}(x_{i}) \leqslant |\lambda_{i}|^{m_{i}+1}\)；于是有 \(|\lambda_{i}|^{m_{i}} < |\lambda_{i}|^{-2}p_{i}(x_{i})\)，由此并根据 (6)，即得关系式 (5)，其中

\[
a = \beta (| \lambda_ {1} | \cdot | \lambda_ {2} | \dots | \lambda_ {n} |) ^ {- 2}. \qquad \text{Q.E.D.}
\]

推论. --- 集 H 是等度连续的，当且仅当：对每个半范数 \(q \in \Gamma\)，存在 \(\prod_{i=1}^{n} E_{i}\) 中 0 的一个邻域，使得诸函数 \(q \circ u\)（\(u \in H\)）在其中一致有界。

条件显然是必要的，而且 prop. 4 的证明表明它蕴含对所有 \(u \in H\) 成立的形如 (5) 的不等式，从而蕴含 H 的等度连续性。

我们把 prop. 4 关于线性映射的特例明确叙述出来。

命题 5. --- 设 E、F 是非离散赋值除环 K 上的两个拓扑向量空间；假设 E（相应地 F）的拓扑由半范数集 \(\Gamma\)（相应地 \(\Gamma'\)）定义。设 H 是 E 到 F 中的线性映射组成的一个集。下列条件等价：

\begin{enumerate}
\def\labelenumi{\alph{enumi})}
\item
  H 等度连续。
\item
  对每个半范数 \(q \in \Gamma'\)，存在半范数的一个有限族 \((p_i)_{1 \leq i \leq n}\)（属于 \(\Gamma\)）及一个数 \(a > 0\)，使得对所有 \(x \in E\) 及所有 \(u \in H\)，
\end{enumerate}

\[
q (u (x)) \leqslant a. \sup _ {1 \leqslant i \leqslant n} p _ {i} (x).\tag{7}
\]

\begin{enumerate}
\def\labelenumi{\alph{enumi})}
\setcounter{enumi}{2}
\tightlist
\item
  对每个半范数 \(q \in \Gamma'\)，映射 \(\sup_{u \in H} (q \circ u)\) 是 \(E\) 上的一个连续半范数。
\end{enumerate}

推论 1. --- 设 \(\mathcal{T}, \mathcal{T}'\) 是向量空间 \(\mathbf{E}\)（在 \(\mathbf{K}\) 上）上分别由两个半范数集 \(\Gamma\) 与 \(\Gamma'\) 定义的拓扑。\(\mathcal{T}\) 比 \(\mathcal{T}'\) 更细，当且仅当：对每个半范数 \(q \in \Gamma'\)，存在半范数的一个有限族 \((p_i)_{1 \leqslant i \leqslant n}\)（属于 \(\Gamma\)）及一个数 \(a > 0\)，使得对所有 \(x \in \mathbf{E}\)，有 \(q(x) \leqslant a . \sup_{1 \leqslant i \leqslant n} p_i(x)\)。

事实上，这表明赋有拓扑 T 的 E 到赋有拓扑 \(T'\) 的 E 上的恒等映射连续。

推论 2. --- 假设 \(\mathcal{T}\) 是拓扑向量空间 \(E\)（在 \(K\) 上的）的拓扑，由一个有向半范数集 \(\Gamma\) 定义；对每个半范数 \(p \in \Gamma\)，设 \(E_p\) 是由 \(E\) 赋有由 \(p\) 定义的拓扑而得的空间。集 \(E'\)（由 \(E\) 上关于 \(\mathcal{T}\) 连续的线性型组成）是诸集 \(E_p'\) 的并，其中 \(E_p'\) 是 \(E_p\) 中连续线性型组成的集（\(p \in \Gamma\)）。

\section{凸集}

\subsection{凸集的定义}

对仿射空间 E 中任意两点 x、y，点 \(\lambda x + \mu y\)（其中 \(\lambda \geqslant 0\)、\(\mu \geqslant 0\)、\(\lambda + \mu = 1\)）组成的集，称为以 x 与 y 为端点的闭线段；当 x = y 时它退化为一个点。该线段中对 x 取补所得的集，称为以 x、y 为端点且在 x 处开、在 y 处闭的线段；当 x = y 时它是空集。最后，以 x、y 为端点的闭线段中对 \(\{x, y\}\) 取补所得的集，称为以 x、y 为端点的开线段；当 x = y 时它是空集。

定义 1. --- 仿射空间 E 的子集 A 称为凸的，如果对 A 的任意两点 x、y，以 x、y 为端点的闭线段都包含在 A 中。

由于 \((1 - \lambda)a + \lambda x = a + \lambda(x - a)\)，这个定义等价于下列定义：集 A 是凸的，如果对每点 \(a \in A\)，A 在以 a 为中心、以 \(\lambda\)（\(0 < \lambda < 1\)）为比的位似变换下的像包含在 A 中（换言之，A 对这些位似变换是稳定的）。

例. --- 1) E 的每个线性仿射流形（特别地，空集）都是凸的。

\begin{enumerate}
\def\labelenumi{\arabic{enumi})}
\setcounter{enumi}{1}
\item
  \(\mathbf{R}\) 中仅有的非空凸集是区间（GT, IV, § 2.4, prop. 1）。
\item
  设 E 是向量空间，\(\|x\|\) 是 E 上的一个范数；由满足 \(\| x\| \leqslant 1\) 的点 x 组成的单位球 B 是凸的，因为关系式 \(\| x\| \leqslant 1\)、\(\| y\| \leqslant 1\) 蕴含：当 \(0\leqslant \lambda \leqslant 1\) 时
\end{enumerate}

\[
\left\| \lambda x + (1 - \lambda) y \right\| \leqslant \lambda \| x \| + (1 - \lambda) \| y \| \leqslant \lambda + (1 - \lambda) = 1.
\]

注记. --- 设 A 是向量空间 E 的一个凸子集；对任意标量 \(\alpha > 0\) 与 \(\beta > 0\)，我们有 \(\alpha A + \beta A = (\alpha + \beta) A\)。换言之，对任意 \(x \in A\)、\(y \in A\)，存在 \(z \in A\) 使得 \((\alpha + \beta) z = \alpha x + \beta y\)；事实上，这个关系式可写成 \(z = \frac{\alpha}{\alpha + \beta} x + \frac{\beta}{\alpha + \beta} y\)，而且我们有 \(\frac{\alpha}{\alpha + \beta} > 0\)、\(\frac{\beta}{\alpha + \beta} > 0\) 及 \(\frac{\alpha}{\alpha + \beta} + \frac{\beta}{\alpha + \beta} = 1\)，由此并利用定义 1 即得结论。

命题 1. --- 设 \((x_{\iota})\) 是凸子集 A 中点的一个族；每个重心 \(\sum_{\iota} \lambda_{\iota} x_{\iota}\)（诸 \(x_{\iota}\) 以正质量 \(\lambda_{\iota}\) 构成的，满足 \(\sum_{\iota} \lambda_{\iota} = 1\)，且除有限多个指标外 \(\lambda_{\iota} = 0\)，cf.~A, II, §9.3）都属于 A。

显然只需考虑指标为 1, 2, \ldots, p 且对每个 i 有 \(\lambda_{i} > 0\) 的情形；当 p = 1 时命题是平凡的；我们对 p 作归纳证明。置 \(\mu = \sum_{i=1}^{p-1} \lambda_{i} > 0\)，并置 \(y = \sum_{i=1}^{p-1} \frac{\lambda_{i}}{\mu} x_{i}\)；归纳假设蕴含 \(y \in A\)。现在由于 \(\lambda_{p} = 1 - \mu\) 且 \(\sum_{i=1}^{p} \lambda_{i} x_{i} = \mu y + (1 - \mu) x_{p}\)，由定义 1 得 \(\sum_{i=1}^{p} \lambda_{i} x_{i}\) 属于 A。

命题 2. --- 设 E 与 F 是两个仿射空间，f 是 E 到 F 中的一个仿射线性映射；则 E 的凸子集在 f 下的像，以及 F 的凸子集在 f 下的原像，都是凸的。

以 x、y 为端点的闭线段在 f 下的像是以 \(f(x)\)、\(f(y)\) 为端点的闭线段，由此得第一个结论。我们推出，F 的闭线段在 f 下的原像包含端点属于它的每个闭线段；由此得 prop. 2 的第二个结论。

特别地，凸集在位似变换或平移下的像是凸集。

命题 3. --- 在仿射空间 E 中，设 H 是由关系式 \(g(x) = 0\) 定义的超平面，其中 g 是 E 上的一个非常数仿射函数。则由关系式 \(g(x) \geqslant 0\)、\(g(x) \leqslant 0\)、\(g(x) > 0\)、\(g(x) < 0\) 定义的诸半空间都是凸的。

因为它们都是 R 的区间在 g 下的原像，从而都是凸的。

采用 prop. 3 的记号，仿射空间中子集 M 的各个点称为位于超平面 H 的同侧（相应地，严格同侧），如果 M 包含在由 \(g(x) \geqslant 0\)、\(g(x) \leqslant 0\)（相应地，\(g(x) > 0\) 或 \(g(x) < 0\)）定义的半空间之一中。

命题 4. --- 仿射空间 E 的凸子集 A 的各个点严格位于超平面 H 的同侧，当且仅当 A 与 H 不相交。

条件显然是必要的。反之，设条件满足，并设 \(g(x) = 0\) 是定义 H 的一个方程（g 是 E 到 R 中的一个仿射线性映射）。集 \(g(\mathbf{A})\) 在 R 中是凸的，因此是一个区间，且 \(0 \notin g(\mathbf{A})\)。故 \(g(x)\) 对所有 \(x \in A\) 具有固定的符号。

\subsection{凸集的交。凸集的乘积}

命题 5. --- 仿射空间 E 的凸子集的任意一族之交是凸的。

本命题由 II, 第 7 页 的定义 1 立即得出。

命题 6. --- 设 \((\mathrm{E}_{\iota})_{\iota\in\mathbf{I}}\) 是向量空间的一个族，且对每个 \(\iota\in I\)，设 \(A_{\iota}\) 是 \(E_{\iota}\) 的一个非空子集。则集 \(A=\prod_{\iota\in I}A_{\iota}\) 在 \(E=\prod_{\iota\in I}E_{\iota}\) 中是凸的，当且仅当对所有 \(\iota\in I\)，集 \(A_{\iota}\) 在 \(E_{\iota}\) 中是凸的。

事实上，每个投影 \(\mathrm{pr}_{\iota}\) 都是线性映射，且我们有 \(A_{\iota} = \mathrm{pr}_{\iota}A\) 与 \(A = \bigcap_{\iota \in I} \mathrm{pr}_{\iota}^{-1}(A_{\iota})\)；本命题由上面的 prop. 2 与 prop. 5 得出。

推论. --- 在空间 \(\mathbf{R}^n\) 中，每个超平行体 (GT, VI, § 1.3) 都是凸子集。

因为它是矩形平行多面体在一个仿射线性映射下的像，而后者由 prop. 6 是凸的。

命题 7. --- 设 A 与 B 是向量空间 E 的两个凸子集。对任意实数 \(\alpha\)、\(\beta\)，集 \(\alpha A + \beta B\)（形如 \(\alpha x + \beta y\) 的点组成的集，其中 x 在 A 中变化，y 在 B 中变化）是凸的。

因为 \(\alpha A + \beta B\) 是凸子集 \(A \times B\)（\(E \times E\) 中的）在线性映射 \((x, y) \mapsto \alpha x + \beta y\)（\(E \times E\) 到 E 中的）下的像。

\subsection{集的凸包}

定义 2. --- 给定仿射空间 E 的子集 A，我们称包含 A 的所有凸集之交为 A 的凸包，也就是说（II, 第 9 页, prop. 5），它是包含 A 的最小凸集。

命题 8. --- 对仿射空间 E 的凸子集的任意族 \((\mathbf{A}_{\iota})_{\iota \in \mathbf{I}}\)，\(\bigcup_{\iota \in \mathbf{I}}\mathbf{A}_{\iota}\) 的凸包恰是线性组合 \(\sum_{\iota \in \mathbf{I}}\lambda_{\iota}x_{\iota}\) 组成的集，其中 \(x_{\iota}\in \mathbf{A}_{\iota},\lambda_{\iota}\geqslant 0\) 对所有 \(\iota \in \mathrm{I}\) 成立（除有限多个指标外 \(\lambda_{\iota} = 0\)），且 \(\sum_{\iota \in \mathbf{I}}\lambda_{\iota} = 1\)。

用 C 表示这些线性组合组成的集，显然 C 包含在包含所有 \(A_{\iota}\) 的每个凸集中（II, 第 8 页, prop. 1）；另一方面，\(A_{\iota} \subset C\)（对每个 \(\iota\)）。剩下只需证明 C 是凸的。设 \(x = \sum_{i} \lambda_{i} x_{i}\)、\(y = \sum_{i} \mu_{i} y_{i}\) 是 C 的两点，\(\alpha\) 是满足 \(0 < \alpha < 1\) 的一个数，写 \(\gamma_{\iota} = \alpha \lambda_{\iota} + (1 - \alpha)\mu_{\iota}\)（对每个 \(\iota \in \mathbf{I}\)），并用 J 表示 I 中使 \(\gamma_{\iota} \neq 0\) 的指标组成的（有限）集。我们可以写 \(\alpha x + (1 - \alpha)y = \sum_{\iota \in \mathbf{J}}\gamma_{\iota}z_{\iota}\)，其中

\[
z _ {\iota} = \gamma_ {\iota} ^ {- 1} (\alpha \lambda_ {\iota} x _ {\iota} + (1 - \alpha) \mu_ {\iota} y _ {\iota})
\]

属于 \(\mathbf{A}_{\iota}\)（对所有 \(\iota \in \mathbf{J}\)）；但 \(\sum_{\iota \in \mathbf{J}} \gamma_{\iota} = \alpha \sum_{\iota \in \mathbf{I}} \lambda_{\iota} + (1 - \alpha) \sum_{\iota \in \mathbf{I}} \mu_{\iota} = 1\)，于是我们看到 \(\alpha x + (1 - \alpha)y \in \mathbf{C}\)。命题得证。

推论 1. --- E 的子集 A 的凸包与线性组合 \(\sum_{i} \lambda_{i} x_{i}\) 组成的集恒同，其中 \((x_{i})\) 是 A 中点的任意有限族，数 \(\lambda_{i} > 0\)（对所有 \(i\)）且 \(\sum_{i} \lambda_{i} = 1\)。

凸集 A 生成的仿射线性流形 (A, II, § 9.3) 的维数称为 A 的维数。

设 E 是向量空间。E 中集 A 的均衡包的凸包 C 称为 A 的均衡凸包（或对称凸包）；显然它是包含 A 的最小对称凸集；它也是 \(A \cup (-A)\) 的凸包，因为 A 的均衡包的每个点都属于以 a 与 -a（其中 \(a \in A\)）为端点的线段。集 C 与线性组合 \(\sum_{i} \lambda_{i} x_{i}\)（其中 \(x_{i} \in A\) 且 \(\sum_{i} |\lambda_{i}| \leqslant 1\)）组成的集重合；因为显然

这个点集是凸的且包含 A 与 -A；只需证明它包含在 C 中，为此只需考虑满足 \(\mu = \sum_{i} |\lambda_{i}| > 0\) 的那些线性组合；于是我们可以写 \(\sum_{i} \lambda_{i} x_{i} = \mu \cdot \sum_{i} \alpha_{i} y_{i}\)，其中 \(\alpha_{i} = \lambda_{i} / \mu\)，\(y_{i} = x_{i}\)（若 \(\lambda_{i} \geqslant 0\)）；而 \(\alpha_{i} = -\lambda_{i} / \mu\)，\(y_{i} = -x_{i}\)（若 \(\lambda_{i} < 0\)）；显然 \(\sum_{i} \alpha_{i} = 1\)，我们的断言得证。

推论 2. --- 设 f 是仿射空间 E 到仿射空间 F 中的一个仿射线性映射；对 E 的每个子集 A，\(f(A)\) 的凸包是 A 的凸包在 f 下的像。

对线性映射与均衡凸包有类似的命题。

\subsection{凸锥}

定义 3. --- 仿射空间 E 的子集 C 称为以 \(x_0\) 为顶点的锥，如果 C 在以 \(x_0\) 为中心、比 \(> 0\) 的所有位似变换下不变。

在本目及下一目中，我们将假设已把所考虑的锥的顶点取为 E 中的原点；即假设 E 是向量空间，并且当我们谈到锥时，应理解该锥以 0 为顶点。形如 \(\lambda a\)（其中 \(\lambda > 0\)，相应地 \(\lambda \geqslant 0\)）的点组成的集（其中 \(a\) 是非零向量）称为以 0 为起点的开半直线（相应地闭半直线）。

以 0 为顶点的锥 C 若满足 \(0 \in C\) 则称为尖锥，否则称为无尖锥。尖锥或者是单点集 \(\{0\}\)，或者是以 0 为起点的一族闭半直线之并。无尖锥是以 0 为起点的开半直线之并（可能为空）。若 C 是无尖锥，则 \(C \cup \{0\}\) 是尖锥。若 C 是尖锥，则 \(C - \{0\}\) 是无尖锥。

若 C 是无尖凸锥，则 \(C \cup \{0\}\) 是尖凸锥。然而，若 C 是尖凸锥，\(C - \{0\}\) 不一定是凸的。我们称尖凸锥为真尖锥，如果它不包含任何过 0 的直线。于是

命题 9. --- 尖凸锥 C 是真尖锥，当且仅当在 C 中对 0 取补所得的无尖锥 C' 是凸的。

若 C 包含一条过 0 的直线，则显然 \(C'\) 不是凸的。现在设 C 是真尖锥，并设 x、y 是 \(C'\) 的两点。以 x、y 为端点的闭线段包含在 C 中；若它包含 0，则 \(\lambda x + (1 - \lambda)y = 0\)（对某个 \(\lambda\)，满足 \(0 < \lambda < 1\)），因此 \(x = \mu y\)，其中 \(\mu < 0\)。于是 C 包含过 0 与 x 的直线，与假设矛盾。

命题 10. --- E 的子集 C 是凸锥，当且仅当 \(C + C \subset C\) 且 \(\lambda C \subset C\)（对所有 \(\lambda > 0\)）。

因为条件「\(\lambda C \subset C\)（对所有 \(\lambda > 0\)）」刻画了锥。若 C 是凸的，我们有 \(C + C = \frac{1}{2}C + \frac{1}{2}C = C\)（II, 第 8 页, 注记）。反之，若锥 C 满足 \(C + C \subset C\)，则当 \(0 < \lambda < 1\) 时，\(\lambda C + (1 - \lambda)C = C + C \subset C\)，这表明 C 是凸的。

推论 1. --- 若 C 是非空凸锥，则 C 生成的向量空间是集 C - C（由 x - y（其中 x、y 在 C 中变化）组成的点集）。

因为，若 V = C - C，则 V 非空，我们有 \(\lambda V = V\)（对所有 \(\lambda \neq 0\)），且 \(V + V = C + C - (C + C) \subset C - C = V\)，这表明 V 是向量子空间。最后，每个包含 C 的向量子空间也包含 V。

推论 2. --- 若 C 是尖凸锥，则包含在 C 中的最大向量子空间是集 \(C \cap (-C)\)。

因为，若 \(W = C \cap (-C)\)，则 W 非空且 \(\lambda W = W\)（对所有 \(\lambda \neq 0\)），又有

\[
\mathrm{W} + \mathrm{W} \subset (\mathrm{C} + \mathrm{C}) \cap (- (\mathrm{C} + \mathrm{C})) \subset \mathrm{C} \cap (- \mathrm{C}) = \mathrm{W},
\]

这表明 W 是向量子空间。显然，包含在 C 中的每个向量子空间也都包含在 W 中。

显然，若 f 是 E 到向量空间 F 中的线性映射，则 E 中凸锥 C 的像 \(f(\mathbf{C})\) 是 F 中的凸锥。E 中（以 0 为顶点的）凸锥的每个交是凸锥。对 E 的每个子集 A，包含 A 的诸凸锥（这样的锥存在，E 本身就是一个）之交是包含 A 的最小凸锥；它称为由 A 生成的凸锥。

命题 11. --- 设 \((\mathbf{C}_{\iota})_{\iota \in \mathbf{I}}\) 是 E 中凸锥的一个族；由诸 \(\mathbf{C}_{\iota}\) 之并生成的凸锥与点 \(\sum_{\iota \in \mathbf{J}} x_{\iota}\) 组成的集恒同，其中 \(\mathbf{J}\) 是 I 的任意有限

子集，且 \(x_{\iota} \in \mathbf{C}_{\iota}\)（对所有 \(\iota \in \mathbf{J}\)）。

事实上，显然这样的点组成的集 C 是包含诸 \(C_{\iota}\) 之并的一个凸锥，且它包含在包含该并的任何凸锥中。

推论. --- 对 E 的任意子集 A，由 A 生成的凸锥与线性组合 \(\sum_{i\in J}\lambda_{i}x_{i}\) 组成的集恒同，其中 \((x_{i})_{i\in J}\) 是 A 中点的任意有限非空族，且 \(\lambda_{i}>0\)（对所有 \(i\in J\)）。

只需看到：若一个凸锥包含 A 的一点 \(x \neq 0\)，则它也包含半直线 \(C_{x}\)（由点 \(\lambda x\)（\(\lambda\) 在正数集中变化）组成），并且 \(C_{x}\) 是凸锥。

命题 12. --- 若 A 是 E 中的凸集，则由 A 生成的凸锥与 \(C = \bigcup_{\lambda > 0} \lambda A\) 恒同。

集 C 显然是锥；只需证明 C 是凸的。设 \(\lambda x, \mu y\) 是 C 的两点（\(\lambda > 0, \mu > 0, x \in A, y \in A\)）。设 \(\alpha, \beta\) 是两个 \textgreater{} 0 的数，满足 \(\alpha + \beta = 1\)。则 \(\alpha \lambda x + \beta \mu y = (\alpha \lambda + \beta \mu) z\)，其中 \(z \in A\)，且 \(\alpha \lambda + \beta \mu > 0\)；故 \(\alpha \lambda x + \beta \mu y \in C\)。

注记. --- 1) 在 prop. 12 的假设下，若 \(0 \notin \mathbf{A}\)，则锥 \(C\) 是无尖锥，从而 \(C \cup \{0\}\) 是真尖锥。

\begin{enumerate}
\def\labelenumi{\arabic{enumi})}
\setcounter{enumi}{1}
\tightlist
\item
  设 A 是 E 中任意凸集；考察凸集 \(A_{1} = A \times \{1\}\)（在空间 \(F = E \times R\) 中）以及由 \(A_{1}\) 生成的以 0 为顶点的凸锥 C。Prop. 12 表明，\(A_{1}\) 是 C 与 F 中超平面 \(E \times \{1\}\) 的交。因此，E 中每个凸集都可以看作 F 中一个以 0 为顶点的凸锥与超平面 \(E \times \{1\}\) 的交在 E 上的投影。
\end{enumerate}

\subsection{序向量空间}

向量空间 E 上的一个预序结构，记作 \(x \preccurlyeq y\) 或 \(y \succcurlyeq x\)，称为与 E 的向量空间结构相容，如果它满足下列两条公理：

\((\mathrm{EO}_{\mathrm{I}})\) 若 \(x \preccurlyeq y\)，则 \(x + z \preccurlyeq y + z\)（对所有 \(z \in \mathrm{E}\)）。

\((\mathrm{EO}_{\mathrm{II}})\) 若 \(x \succcurlyeq 0\)，则 \(\lambda x \succcurlyeq 0\)（对每个标量 \(\lambda \geqslant 0\)）。

赋有这两个结构的向量空间 E 称为预序向量空间（相应地，当 E 上的预序关系是一个序时，称为序向量空间）。

注意，公理 \((EO_{I})\) 意味着 E 的预序结构与加法群结构相容，也就是说，赋有这两个结构的 E 是一个预序群 (A, VI, 第 3 页)。

例. --- 在定义在 A 上的所有实值函数组成的空间 \(E = R^{A}\) 上，由「对所有 \(t \in A\)，\(x(t) \leqslant y(t)\)」给出的序关系与 E 的向量空间结构相容。

命题 13. --- (i) 预序向量空间 E 中由 \(\succcurlyeq 0\) 的元素组成的集 P 是尖凸锥。

\begin{enumerate}
\def\labelenumi{(\roman{enumi})}
\setcounter{enumi}{1}
\item
  反之，若 \(\mathbf{P}\) 是 \(\mathbf{E}\) 中的一个尖凸锥，则关系 \(y - x \in \mathbf{P}\) 是 \(\mathbf{E}\) 上的一个预序关系，且由它定义的预序结构是唯一与 E 的向量空间结构相容且使 P 为由 \(\succcurlyeq 0\) 的元素组成之集的预序结构。
\item
  关系 \(y - x \in \mathbf{P}\)（其中 \(\mathbf{P}\) 是尖凸锥）是 \(\mathbf{E}\) 上的序关系，当且仅当 \(\mathbf{P}\) 是真尖锥。
\item
  公理 \((\mathrm{EO}_{\mathrm{I}})\) 与 \((\mathrm{EO}_{\mathrm{II}})\) 蕴含 \(P + P \subset P\) 及 \(\lambda P \subset P\)（对所有 \(\lambda > 0\)）。由于 \(0 \in P\)，由此推出 P 是尖凸锥（II, 第 11 页, prop. 10）。
\item
  反之，若 \(\mathbf{P}\) 是尖凸锥，则关系 \(\mathbf{P} + \mathbf{P} \subset \mathbf{P}\) 蕴含关系 \(y - x \in \mathbf{P}\) 是与 E 的加法群结构相容的预序 (A, VI, 第 3 页, prop. 3)；显然，把它写成 \(x \preccurlyeq y\) 时，集 \(\mathbf{P}\) 与由 \(x \geqslant 0\) 组成的集恒同；进而，关系「\(\lambda \mathbf{P} \subset \mathbf{P}\)（对所有 \(\lambda \geqslant 0\)）」表明公理 \((\mathrm{EO}_{\mathrm{II}})\) 满足。(iii) 说 \(\mathbf{P}\) 是真尖锥，意味着 \(\mathbf{P} \cap (-\mathbf{P}) = \{0\}\)（II, 第 11 页, cor. 2），从而 \(y - x \in \mathbf{P}\) 是序关系。
\end{enumerate}

例. --- * 设 H 是一个实 Hilbert 空间；在 H 的连续自同态组成的向量空间 \(\mathcal{L}(\mathrm{H})\) 中，正埃尔米特自同态组成一个真尖凸锥；因此这个锥定义了 \(\mathcal{L}(\mathrm{H})\) 上一个与其向量空间结构相容的序结构，且其中关系 \(A \leqslant B\) 意指 \(B - A\) 是正埃尔米特自同态。*

对向量空间 E 中任意尖凸锥 P，集 \(P \cap (-P)\) 是 E 的一个向量子空间 H（II, 第 11 页, cor. 2）。P 在 E/H 中的典范像 \(P'\) 是凸锥，且 \(P'\) 在 E 中的原像是 P。于是 \(P' \cap (-P') = \{0\}\)，从而 \(P'\) 在 E/H 上定义一个与其向量空间结构相容的序结构。

预序向量空间 E 上的线性型 f 称为正的，如果 E 中 \(x \geqslant 0\) 蕴含 \(f(x) \geqslant 0\)。或者等价地说，如果 E 中由 \(\geqslant 0\) 的元素组成的凸锥 P 包含在由满足 \(f(x) \geqslant 0\) 的 x 组成的半空间中。显然，在 E 的对偶 E* 中，正线性型组成的集是一个尖凸锥。

\subsection{拓扑向量空间中的凸锥}

命题 14. --- 在拓扑向量空间 E 中，凸集（相应地凸锥）的闭包是凸集（相应地具有相同顶点的凸锥）。

因为，设 A 是凸集；映射 \((x, y) \mapsto \lambda x + (1 - \lambda)y\)（其中 \(0 < \lambda < 1\)）在 E × E 中连续且把 A × A 映入 A；于是（GT, I, §2.1, th. 1）它把 \(\overline{A} \times \overline{A}\) 映入 \(\overline{A}\)，这表明 \(\overline{A}\) 是凸的。类似地，若 C 是以 0 为顶点的凸锥，则 \(\overline{C} + \overline{C} \subset \overline{C}\) 且 \(\lambda\overline{C} \subset \overline{C}\)（对所有 \(\lambda > 0\)）。

定义 4. --- 对拓扑向量空间 E 中任意集 A，包含 A 的所有闭凸集之交称为 A 的闭凸包；它是包含 A 的最小闭凸集。

由 prop. 14，A 的闭凸包是 A 的凸包的闭包；显然它与 \(\overline{A}\) 的闭凸包相同。

类似地，我们称包含 A 的最小对称闭凸集为 A 的对称闭凸包（或均衡闭凸包）；它是 A 的对称凸包的闭包（II, 第 10 页）；它也是 \(\overline{A}\) 的对称闭凸包。

命题 15. --- 设 \(A_{i}\)（\(1 \leqslant i \leqslant n\)）是 Hausdorff 拓扑向量空间 E 中有限个紧凸集。则诸 \(A_{i}\) 之并的凸包是紧的（因此与该并的闭凸包相同）。

设 B 是 \(R^{n}\) 中由点 \((\lambda_{1},\lambda_{2},\ldots,\lambda_{n})\)（其中 \(\lambda_{i}\geqslant0(1\leqslant i\leqslant n)\) 且 \(\sum_{i=1}^{n}\lambda_{i}=1\)）定义的紧集。定义 \(B\times\prod_{i=1}^{n}A_{i}\subset R^{n}\times E^{n}\) 到 E 中的一个连续映射（由公式 \((\lambda_{1},\lambda_{2},\ldots,\lambda_{n},x_{1},x_{2},\ldots,x_{n})\mapsto\sum_{i=1}^{n}\lambda_{i}x_{i}\)）。\(\bigcup_{i=1}^{n}A_{i}\) 的凸包 C 是 \(B\times\prod_{i=1}^{n}A_{i}\) 在该映射下的像；由于 \(B\times\prod_{i=1}^{n}A_{i}\) 是紧的且 E 是 Hausdorff 的，由此推出 C 是紧的。

推论 1. --- 在 Hausdorff 拓扑向量空间中，有限集的凸包是紧的。

推论 2. --- 在拓扑向量空间 E 中，有限集的凸包是预紧的。

事实上，设 \(j\) 是 \(\mathbf{E}\) 到其 Hausdorff 完备化 \(\hat{\mathbf{E}}\) 中的典范映射；若 \(\mathbf{C}\) 是 \(\mathbf{A}\) 的凸包，则 \(j(\mathbf{C})\) 是有限集 \(j(\mathbf{A})\)（在 \(\hat{\mathbf{E}}\) 中的）的凸包，故 \(j(\mathbf{C})\) 是紧的（cor. 1），从而 \(\mathbf{C}\) 是预紧的（GT, II, § 4.2）。

命题 16. --- 设 A 是拓扑向量空间 E 的一个凸子集，且至少有一个内点 \(x_{0}\)。则对任意点 \(x \in \overline{A}\)，以 \(x_{0}\)、x 为端点的开线段的每个点都属于 A 的内部。

设 y 是该线段中任意一点，f 是以 y 为中心、\(\lambda < 0\) 为比且把 \(x_{0}\) 变为 x 的位似变换。若 V 是包含在 A 中的 \(x_{0}\) 的一个开邻域，则 \(f(\mathrm{V})\) 是 x 的邻域，因此包含一点 \(f(z) \in \mathrm{A}\)；现在

\[
f (z) - y = \lambda (z - y) = \lambda (z - f (z)) + \lambda (f (z) - y),
\]

于是 \(y - f(z) = \frac{\lambda}{\lambda - 1} (z - f(z))\)，从而 \(y\) 被变为 \(z\)（由位似变换 \(g\)，其中心为 \(f(z)\)，比为 \(\mu = \lambda / (\lambda - 1)\)）；由于 \(0 < \mu < 1\)，\(g\) 把 \(V\) 变为包含在 \(A\) 中的 0 的邻域。命题得证。

推论 1. --- 内部 \(\mathring{\mathrm{A}}\)（凸集 \(\mathrm{A}\) 的）本身是凸集；若 \(\mathring{\mathrm{A}}\) 非空，则它与 \(\overline{\mathrm{A}}\) 的内部重合，且 \(\overline{\mathrm{A}}\) 是与 \(\mathring{\mathrm{A}}\) 的闭包重合的凸集。

由 prop. 16 推出：若 \(\mathring{\mathrm{A}}\) 非空，则它是凸集且 \(\overline{\mathrm{A}}\) 的每点都是 \(\mathring{\mathrm{A}}\) 的接触点。其次我们证明 \(\overline{\mathrm{A}}\) 的每个内点属于 \(\mathring{\mathrm{A}}\)。设 \(x\) 是 \(\overline{\mathrm{A}}\) 的一个内点，并不妨设 \(x = 0\)。设 \(\mathrm{V}\) 是包含在 \(\overline{\mathrm{A}}\) 中的 0 的一个对称邻域，并设 \(y \in \mathring{\mathrm{A}} \cap \mathrm{V}\)；于是 \(-y \in \overline{\mathrm{A}}\)，因此由 prop. 16，我们看到 \(0 \in \mathring{\mathrm{A}}\)（若 \(y \neq 0\)）；当 \(y = 0\) 时这显然成立。

推论 2. --- 内部 \(\mathring{\mathrm{C}}\)（凸锥 \(\mathrm{C}\) 的）本身是凸锥；若 \(\mathring{\mathrm{C}}\) 非空，则它与 \(\overline{\mathrm{C}}\) 的内部重合，且 \(\overline{\mathrm{C}}\) 是与 \(\mathring{\mathrm{C}}\) 的闭包重合的尖凸锥。

由于以 0 为中心、\textgreater{} 0 为比的位似变换把 C 变为自身，它们对 \(\overset{\circ}{C}\) 也如此，故 \(\overset{\circ}{C}\) 是锥；推论的其余部分由 cor. 1 及显然的注记（若 C 非空，则 \(\overline{C}\) 包含 C 的顶点）得出。

设 H 是 R 上拓扑向量空间 E 中的一个闭超平面；它有一个形如 \(f(x) = \alpha\) 的方程，其中 f 是 E 上不恒为零的连续线性型（I, 第 13 页, th. 1）。于是由 \(f(x) \leqslant \alpha\) 与 \(f(x) \geqslant \alpha\) 分别定义的闭半空间是闭凸集；它们的补集，即由 \(f(x) > \alpha\) 与 \(f(x) < \alpha\) 分别定义的集，是开凸集。我们称这些半空间是由 H 确定的闭（相应地开）半空间。

命题 17. --- 在拓扑向量空间 E 中，设 A 是至少有一个内点的集，且其所有点都位于某个超平面 H 的同侧。则 H 是闭的，A 的内点严格位于 H 的同侧，而 A 的接触点位于 H 的同侧。特别地，开（相应地闭）半空间由闭超平面确定。

事实上，设 H 包含原点且 \(f(x) = 0\) 是 H 的一个方程；不妨设 \(f(x) \geqslant 0\)（对所有 \(x \in A\)）。由满足 \(f(y) > -1\) 的点 y 组成的半空间至少包含一个内点，并且通过平移，由满足 \(f(y) > 0\) 的点 y 组成的半空间也是如此；这表明 H 是闭的（I, 第 11 页, corollary）。于是我们知道 f 是 E 到 R 上的一个严格态射（I, 第 13 页, corollary），因此 \(f(\mathring{\mathrm{A}})\) 是 R 中的开集。这个集不能包含 0，否则它将包含 \textless{} 0 的数，与假设矛盾；从而它包含在开区间 \(\bigg]0, +\infty[\) 中。另一方面，由满足 \(f(y) \geqslant 0\) 的 y 组成的半空间是闭的且包含 A，因此它包含 \(\overline{A}\)。

推论. --- 设 P 是拓扑向量空间 E 中至少有一个内点的尖凸锥。则 E 上每个不恒为零且关于由 P 定义的预序结构（II, 第 13 页）为正的线性型 f 必定连续。进而，若 x 是 P 的内点，则 \(f(x) > 0\)；若 x 是 P 的接触点，则 \(f(x) \geqslant 0\)。

把 prop. 17 应用于 \(\mathbf{A} = \mathbf{P}\) 的情形，其中 \(\mathbf{H}\) 是以 \(f(x) = 0\) 为方程的超平面。

注记. --- 在拓扑向量空间 E 中，每个凸集 C 都是连通的。事实上，若 \(a \in C\)，则 C 是以 a 为一个端点且在 a 处闭的一些线段之并；这些线段都是连通的，结论由 GT, I, § 11.1, prop. 2 得出。

\subsection{序向量空间上的拓扑}

设 E 是序向量空间。E 上的拓扑称为与 E 的序向量空间结构相容，如果它既与 E 的向量空间结构相容，又满足下列公理：

(TO) 诸 \(x\)（\(x \geqslant 0\)）组成的凸锥在 E 中是闭的。

赋有相容拓扑的序向量空间 E 称为序拓扑向量空间。

例. --- 赋有通常拓扑及其因子的序结构之乘积这一序结构的空间 \(R^{n}\) 是序拓扑向量空间。另一方面，当 \(n \geqslant 2\) 且 \(R^{n}\) 赋有字典序 (S, III, § 2.6) 时，通常拓扑与 \(R^{n}\) 的序向量空间结构不相容。

设 \(\mathbf{A}\) 是一个集；向量空间 \(\mathcal{B}(\mathbf{A};\mathbf{R})\)（由定义在 \(\mathbf{A}\) 上的实值有界函数组成），赋有由范数 \(\| x\| = \sup_{t\in \mathbf{A}}|x(t)|\) 定义的拓扑以及由

\(R^{A}\) 的乘积序结构诱导的序结构，是一个序拓扑向量空间。

在序拓扑向量空间 E 中，由元素 \(x \leqslant 0\) 组成的集是闭的；由于平移是同胚，我们推出，对所有 \(a \in E\)，由元素 \(x \geqslant a\)（相应地 \(x \leqslant a\)）组成的集是闭的。由于 \(\{0\}\) 是集 \(x \geqslant 0\) 与 \(x \leqslant 0\) 之交，由此推出 \(\{0\}\) 是闭的，从而 E 是 Hausdorff 的。

命题 18. --- 在序拓扑向量空间 E 中，设 H 是由关系 \(\leqslant\) 定向的一个集。若 H 的截段滤子在 E 中有极限，则该极限是 H 的上确界。

因为，设 \(b = \lim_{x\in \mathbf{H}}x\)；对每个 \(y\in \mathbf{H}\)，诸 \(x\in \mathbf{H}\)（满足 \(x\geqslant y\)）组成的集是

H 的截段滤子中的一个集，因此 b 是该集的一个接触点；但由于集 \(x \geqslant y\) 在 E 中是闭的，我们有 \(b \geqslant y\)，故 b 是 H 的一个上界。另一方面，若 a 是 H 的一个上界，则 H 包含在闭集 \(x \leqslant a\) 中；由于 b 是 H 的一个接触点，我们有 \(b \leqslant a\)，证毕（II, 第 72 页, exerc. 42）。

\subsection{凸函数}

定义 5. --- 设 X 是仿射空间 E 的一个凸子集。定义在 X 上的实值有限函数称为凸的（相应地严格凸的），如果对 X 中任意两个不同的点 x、y 及任意实数 \(\lambda\)（\(0 < \lambda < 1\)），我们有：

\[
f (\lambda x + (1 - \lambda) y) \leqslant \lambda f (x) + (1 - \lambda) f (y)\tag{1}
\]

（相应地

\[
f (\lambda x + (1 - \lambda) y) <   \lambda f (x) + (1 - \lambda) f (y)).\tag{2}
\]

当 E = R 时，凸函数的这一定义与 FVR, I, 第 32 页 中的定义相同。进而，f 在 X 中是凸的（相应地严格凸的），当且仅当对每条仿射直线 \(D \subset E\)，f 在 \(X \cap D\) 上的限制在 \(X \cap D\) 中是凸的（相应地严格凸的）。

例. --- 若 \(f\) 是 \(\mathbf{E}\) 上的一个仿射线性函数，则 \(f\) 与 \(f^2\) 都是 \(\mathbf{E}\) 上的凸函数；对 \(f\) 而言这是显然的，因为

\[
f (\lambda x + (1 - \lambda) y) = \lambda f (x) + (1 - \lambda) f (y);
\]

另一方面，若 \(\alpha = f(x)\)、\(\beta = f(y)\)，则；

\[
\lambda \alpha^ {2} + (1 - \lambda) \beta^ {2} - (\lambda \alpha + (1 - \lambda) \beta) ^ {2} = \lambda (1 - \lambda) (\alpha - \beta) ^ {2} \geqslant 0
\]

其中 \(0 < \lambda < 1\)；进而，\(f^2\) 在仿射直线 \(\mathbf{D} \subset \mathbf{E}\) 上的限制是严格凸的（若 \(f|\mathbf{D}\) 不是常数）。

实值函数 \(f\)（定义在 \(\mathbf{X}\) 上的）称为凹的（相应地严格凹的），如果 \(-f\) 是凸的（相应地严格凸的）。也就是说，对任意两个不同的点 \(x, y\)（\(\mathbf{X}\) 中的）及每个数 \(\lambda\)（满足 \(0 < \lambda < 1\)），我们有

\[
f (\lambda x + (1 - \lambda) y) \geqslant \lambda f (x) + (1 - \lambda) f (y)
\]

（相应地

\[
f (\lambda x + (1 - \lambda) y) > \lambda f (x) + (1 - \lambda) f (y)).
\]

X 到 R 中的映射称为仿射的，如果它既是凸的又是凹的（cf.~II, 第 78 页, exerc. 11）。

命题 19. --- 设 X 是仿射空间 E 中的凸集；f 是定义在 X 上的实值函数。用 F（相应地 F'）表示点 \((x, a) \in \mathbf{E} \times \mathbf{R}\)（满足 \(x \in X\) 且 \(f(x) \leqslant a\)（相应地 \(x \in X\) 且 \(f(x) < a\)））组成的集。则下列条件等价：

\begin{enumerate}
\def\labelenumi{\alph{enumi})}
\item
  函数 \(f\) 是凸的。
\item
  集 \(\mathbf{F}\) 在仿射空间 \(\mathbf{E} \times \mathbf{R}\) 中是凸的。
\item
  集 \(\mathbf{F}'\) 在仿射空间 \(\mathbf{E} \times \mathbf{R}\) 中是凸的。
\end{enumerate}

我们证明 \(a) \Rightarrow c)\)。设 \((x, a)\) 与 \((y, b)\) 是 \(F'\) 的两点且 \(0 < \lambda < 1\)，则 \(f(x) < a\)，\(f(y) < b\)，而若 \(f\) 是凸的，则

\[
f (\lambda x + (1 - \lambda) y) \leqslant \lambda f (x) + (1 - \lambda) f (y) <   \lambda a + (1 - \lambda) b
\]

这表明点 \(\lambda(x, a) + (1 - \lambda)(y, b)\)（\(\mathbf{E} \times \mathbf{R}\) 中的）属于 \(\mathbf{F}'\)。于是 \(\mathbf{F}'\) 是凸的。

其次我们证明 \(c) \Rightarrow b)\)。若 \((x, a), (y, b)\) 是 \(F\) 的两点，则对每个 \(\varepsilon > 0\)，\((x, a + \varepsilon)\) 与 \((y, b + \varepsilon)\) 属于 \(F'\)，且当 \(0 < \lambda < 1\) 时，\((\lambda x + (1 - \lambda)y, \lambda a + (1 - \lambda)b + \varepsilon)\) 亦然；由 \(F\) 的定义，这就蕴含 \((\lambda x + (1 - \lambda)y, \lambda a + (1 - \lambda)b)\) 属于 \(F\)。

最后 \(b) \Rightarrow a)\)，因为（采用上面的记号）若 \((\lambda x + (1 - \lambda)y, \lambda a + (1 - \lambda)b)\) 属于 F，则

\[
f (\lambda x + (1 - \lambda) y) \leqslant \lambda a + (1 - \lambda) b
\]

只要 \(a \geqslant f(x)\) 且 \(b \geqslant f(y)\)；由此得出 (1)，从而 \(f\) 是凸的。

推论. --- 若 \(f\) 在 \(\mathbf{X}\) 中是凸的，则对所有 \(\alpha \in \mathbf{R}\)，诸 \(x \in \mathbf{X}\)（满足 \(f(x) \leqslant \alpha\)（相应地 \(f(x) < \alpha\)））组成的集是凸的。

事实上，它是 F（相应地 F'）与超平面 \(E \times \{\alpha\}\)（\(E \times R\) 中的）之交在 E 上的投影。

命题 20. --- 设 f 是定义在仿射空间 E 的凸集 X 上的凸函数。则对任意族 \((x_{i})_{1\leqslant i\leqslant p}\)（X 的 \(p\) 个点的）及任意族 \((\lambda_{i})_{1\leqslant i\leqslant p}\)（\(p\) 个实数的，均 \(\geqslant 0\)，满足 \(\sum_{i = 1}^{p}\lambda_{i} = 1\)），我们有：

\[
f \left(\sum_ {i = 1} ^ {p} \lambda_ {i} x _ {i}\right) \leqslant \sum_ {i = 1} ^ {p} \lambda_ {i} f \left(x _ {i}\right).\tag{3}
\]

若 \(f\) 是严格凸的且 \(\lambda_i > 0\)（对所有 \(i\)），则

\[
f \left(\sum_ {i = 1} ^ {p} \lambda_ {i} x _ {i}\right) <   \sum_ {i = 1} ^ {p} \lambda_ {i} f \left(x _ {i}\right),\tag{4}
\]

除非诸 \(x_{i}\) 全相等。

不等式 (3) 由上面的 II, prop. 19 及 II, 第 8 页, prop. 1 得出。设诸 \(x_{i}\) 不全相等（这蕴含 \(p \geqslant 2\)）且诸 \(\lambda_{i}\) 全 \(>0\)；则点 \(z = \sum_{i=1}^{p} \lambda_{i} x_{i}\) 至少与某个 \(x_{i}\) 不同。不妨设 \(z \neq x_{1}\)，写 \(z = \lambda_{1} x_{1} + (1 - \lambda_{1}) y_{1}\)，其中 \(y_{1} = \sum_{i=2}^{p} \frac{\lambda_{i}}{1 - \lambda_{1}} x_{i}\)。则 \(y_{1} \neq x_{1}\)，且由于 \(0 < \lambda_{1} < 1\)，由假设我们有

\[
f (z) <   \lambda_ {1} f (x _ {1}) + (1 - \lambda_ {1}) f (y _ {1}).
\]

但由 (3) 有 \(f(y_{1}) \leqslant \sum_{i=2}^{p} \frac{\lambda_{i}}{1 - \lambda_{1}} f(x_{i})\)，由此即得不等式 (4)。

\subsection{凸函数的运算}

设 \(\mathbf{X}\) 是仿射空间 \(\mathbf{E}\) 中的凸集。若 \(f_{i}(1 \leqslant i \leqslant p)\) 是定义在 \(\mathbf{X}\) 上的有限个凸函数，且 \(c_{i} (1 \leqslant i \leqslant p)\) 是 \(\geqslant 0\) 的数，则函数 \(f = \sum_{i=1}^{p} c_{i} f_{i}\) 在 \(\mathbf{X}\) 上是凸的。

若 \((f_{i})\) 是定义在 \(\mathbf{X}\) 上的凸函数的任意族，且上包络 \(g\)（该族在 \(\mathbf{X}\) 中的）是有限的，则 \(g\) 是凸的。

最后，若 H 是定义在 X 上的凸函数组成的一个集，F 是 H 上在 X 中简单收敛于有限实值函数 \(f_{0}\) 的滤子，则 \(f_{0}\) 在 X 上是凸的。

\subsection{开凸集上的凸函数}

命题 21. --- 设 f 是定义在拓扑向量空间 E 中非空开凸集 X 上的凸函数。则 f 连续，当且仅当它限制在 X 的某个非空开子集 U 上时有上界。

条件显然是必要的，我们证明它是充分的。设 \(x_{0} \in X\) 是使 f 在 \(x_{0}\) 的某邻域 V 中有上界的一点；我们先证明 \(f\) 在 \(x_0\) 处连续。通过平移，我们可以限于 \(x_0 = 0\) 且 \(f(x_0) = 0\) 的情形；此外还可设邻域 \(V\) 是均衡的（I, 第 7 页, prop. 4）。设 \(f(x) \leqslant a\) 在 \(V\) 中成立；对每个 \(\varepsilon, 0 < \varepsilon < 1\)，我们注意到若 \(x \in \varepsilon V\)，则 \(x / \varepsilon \in V\) 且 \(-x / \varepsilon \in V\)。把 II, 第 16 页 的不等式 (1) 应用于点 \(x / \varepsilon\)、0 及数 \(\lambda = \varepsilon\)，得 \(f(x) \leqslant \varepsilon f(x / \varepsilon) \leqslant \varepsilon a\)；把它应用于点 \(x\)、\(-x / \varepsilon\) 及数 \(\lambda = 1 / (1 + \varepsilon)\)，得 \(f(x) \geqslant -\varepsilon f(-x / \varepsilon) \geqslant -\varepsilon a\)。于是 \(f(x)\) 在 \(\varepsilon V\) 中任意小（若 \(\varepsilon\) 充分小），故 \(f\) 在 \(x = 0\) 处连续。

现在设 \(y\) 是 \(X\) 中任意一点；由于 \(X\) 是开的，存在数 \(\rho > 1\) 使得 \(z = \rho y\) 属于 \(X\)。设 \(g\) 是位似变换 \(x \mapsto \lambda x + (1 - \lambda)z\)（以 \(z\) 为中心、\(\lambda = 1 - \frac{1}{\rho}\) 为比，它把 0 变为 \(y\)）；对每点 \(g(x) \in g(V)\)，由 (1) 有

\[
f (g (x)) \leqslant \lambda f (x) + (1 - \lambda) f (z) \leqslant \lambda a + (1 - \lambda) f (z).
\]

于是 \(f\) 在 \(y\) 的一个邻域中有上界，从而由第一部分知其在 \(y\) 处连续。命题得证。

推论. --- 每个凸函数 \(f\)（定义在开凸集 \(X\)（\(\mathbf{R}^n\) 中的）上的）都在 \(X\) 中连续。

我们可以设 X 非空。于是在 X 中存在 \(n + 1\) 个仿射无关的点 \(a_{i}\)（\(0 \leqslant i \leqslant n\)），且这些点的凸包 S 包含由点 \(\sum_{i=0}^{n} \lambda_{i} a_{i}\)（\(0 < \lambda_{i} < 1\)（对所有 i）且 \(\sum_{i=0}^{n} \lambda_{i} = 1\)）组成的非空开集。由 II, 第 17 页, prop. 20，f 在 S 中有上界，因此是连续的。

在无穷维拓扑向量空间中，一般存在不连续的线性型（II, 第 80 页, exerc. 25），从而存在在任何点都不连续的凸函数。

\subsection{半范数与凸集}

设 \(\mathbf{E}\) 是 \(\mathbf{R}\) 上的向量空间；映射 \(p\)（\(\mathbf{E}\) 到 \(\mathbf{R}\) 中的）称为正齐次的，如果对每个 \(\lambda \geqslant 0\) 及所有 \(x \in \mathbf{E}\)，我们有

\[
p (\lambda x) = \lambda p (x).\tag{5}
\]

E 上的正齐次函数 p 是凸的，当且仅当它对 E 的所有 x、y 满足 II, 第 1 页 的公理 \((\mathrm{SN}_{\mathrm{II}})\)：

\[
p (x + y) \leqslant p (x) + p (y).\tag{6}
\]

事实上，若 \(p\) 是凸的，则对 \(x, y\)（\(E\) 中的），

\[
p \left(\frac {1}{2} (x + y)\right) \leqslant \frac {1}{2} p (x) + \frac {1}{2} p (y)
\]

且由 (5)，这个关系式等价于 (6)。反之，若 (6) 成立，则对所有 \(\lambda\)（满足 \(0 < \lambda < 1\)），我们也有

\[
p (\lambda x + (1 - \lambda) y) \leqslant p (\lambda x) + p ((1 - \lambda) y) = \lambda p (x) + (1 - \lambda) p (y)
\]

这里利用了 (5)。

E 上的凸正齐次函数称为次线性函数。

若 \(p\) 是定义在 \(\mathbf{E}\) 上的次线性函数，则由 II, § 2.8, corollary，对所有 \(a > 0\)，集 \(\mathrm{V}(p, a)\)（相应地 \(\mathrm{W}(p, a)\)）——由点 \(x \in \mathrm{E}\)（满足 \(p(x) \leqslant a\)（相应地 \(p(x) < a\)））组成——是凸的；进而这个集是吸收的，因为对所有 \(x \in \mathrm{E}\)，存在 \(\lambda > 0\) 使得 \(p(\lambda x) = \lambda p(x) < a\)。

这个结果有一个部分逆命题：

命题 22. --- 设 A 是向量空间 E 中包含 0 的一个凸集。对所有 \(x \in \mathbf{E}\)，置

\[
p _ {\mathrm{A}} (x) = \inf _ {\rho > 0, x \in \rho \mathrm{A}} \rho\tag{7}
\]

\((0 \leqslant p_{\mathbf{A}}(x) \leqslant \infty)\)。函数 \(p_{\mathbf{A}}\) 满足

\[
p _ {\mathrm{A}} (x + y) \leqslant p _ {\mathrm{A}} (x) + p _ {\mathrm{A}} (y), \quad p _ {\mathrm{A}} (\lambda x) = \lambda p _ {\mathrm{A}} (x)\tag{8}
\]

对所有 \(x, y\)（\(\mathbf{E}\) 中的）及 \(\lambda > 0\) 成立。若用 \(\mathrm{V}(p_{\mathbf{A}}, \alpha)\)（相应地 \(\mathrm{W}(p_{\mathbf{A}}, \alpha)\)）表示诸 \(x \in \mathbf{E}\)（满足 \(p_{\mathbf{A}}(x) \leqslant \alpha\)（相应地 \(p_{\mathbf{A}}(x) < \alpha\)））组成的集，则

\[
\mathrm{W} (p _ {\mathrm{A}}, 1) \subset \mathrm{A} \subset \mathrm{V} (p _ {\mathrm{A}}, 1).\tag{9}
\]

若 \(\mathbf{A}\) 是吸收的，则 \(p_{\mathbf{A}}\) 是有限的（从而是次线性的）。

由于关系式 \(x \in \rho A\) 与 \(\lambda x \in \lambda \rho A\) 当 \(\lambda > 0\) 时等价，我们有 \(p_A(\lambda x) = \lambda p_A(x)\)（对 \(\lambda > 0\)）。设 \(x, y\) 是 \(E\) 的两点。若 \(x\)（相应地 \(y\)）不被 \(A\) 吸收，则 \(p_A(x) = +\infty\)（相应地 \(p_A(y) = +\infty\)），而不等式 \(p_A(x + y) \leqslant p_A(x) + p_A(y)\) 显然成立。设存在 \(\alpha > 0\)、\(\beta > 0\) 使得 \(x \in \alpha A\) 且 \(y \in \beta A\)；则 \(x + y \in \alpha A + \beta A = (\alpha + \beta)A\)（II, 第 8 页, 注记）；从而 \(p_A(x + y) \leqslant p_A(x) + p_A(y)\)。包含关系 \(A \subset V(p_A, 1)\) 显然成立。包含关系 \(W(p_A, 1) \subset A\) 则由 \(A\) 是凸的且包含 0 得出。最后，若 \(A\) 是吸收的，则 \(p_A\) 显然是有限的。

由 (7) 定义的函数 \(p_{A}\) 称为凸集 A 的度规。若 A 是吸收的且对称的，则 \(p_{A}\) 是一个半范数。

命题 23. --- 设 E 是拓扑向量空间。若 A 是包含 0 的开凸集，则 \(p_{A}\) 是有限的且连续的，且 \(\mathbf{A} = \mathbf{W}(p_{\mathbf{A}}, 1)\)。若 A 是包含 0 的闭凸集，则 \(p_{A}\) 是下半连续的，且 \(\mathbf{A} = \mathbf{V}(p_{\mathbf{A}}, 1)\)。

若 \(\mathbf{A}\) 是开的且包含 0，则它是吸收的。对 \(x \in \mathbf{A}\)，存在 \(\rho < 1\) 使得 \(x / \rho \in \mathbf{A}\)，从而 \(p_{\mathbf{A}}(x) < 1\)；这与 (9) 结合便给出 \(\mathbf{A} = \mathbf{W}(p_{\mathbf{A}}, 1)\)。由于凸函数 \(p_{\mathbf{A}}\) 在开集 \(\mathbf{A}\) 中有上界，它在 E 中连续（II, 第 18 页, prop. 21）。

设 A 是闭的且包含 0。对每个 \(x \in E\)（满足 \(p_{\mathbf{A}}(x) \leqslant 1\)），我们有 \(x \in \rho A\)（对所有 \(\rho > 1\)），因此由于 A 是闭的，\(x \in A\)；结合 (9)，这表明 \(\mathbf{A} = \mathrm{V}(p_{\mathbf{A}}, 1)\)。于是对所有 \(\mu > 0\)，\(\mu A\) 是诸 \(x\)（满足 \(p_{\mathbf{A}}(x) \leqslant \mu\)）组成的集；由于在 E 中 \(p_{\mathbf{A}}(x) \geqslant 0\)，这表明 \(p_{\mathbf{A}}\) 在 E 中下半连续（GT, IV, § 6.2）。

正次线性函数 \(p\)（\(\mathbf{E}\) 上的）是每个凸集 \(\mathbf{A}\)（其中 \(\mathrm{W}(p,1)\subset \mathrm{A}\subset \mathrm{V}(p,1)\)）的度规。

\section{Hahn-Banach 定理（解析形式）}

\subsection{正线性型的延拓}

命题 1. --- 设 E 是预序向量空间，V 是 E 的一个向量子空间，使得 E 的每个元素都被 V 的某个元素上有界。给定 V 上关于 V 的预序向量空间结构（由 E 的结构诱导）为正的线性型 f，存在 E 上正线性型组成的一个非空集 \(S_{f}\)，其中每个线性型都是 f 的延拓。若 \(h \in S_{f}\)，则诸值 \(h(a)\)（对 \(a \in E\)）都位于区间 \([\alpha', \alpha'']\) 中，其中

\[
\alpha^ {\prime} = \sup _ {z \in \mathbf {V}, z \leqslant a} f (z), \quad \alpha^ {\prime \prime} = \inf _ {y \in \mathbf {V}, y \geqslant a} f (y).\tag{1}
\]

I. 特殊情形。

首先设 \(E = V + R a\)。由于当 \(a \in V\) 时命题是平凡的，我们限于 \(a \notin V\) 的情形。对 V 的假设蕴含集 \(A''\)（诸 \(y \in V\)（满足 \(a \leqslant y\)）组成的）非空；类似地，集 \(A'\)（诸 \(z \in V\)（满足 \(-z \geqslant -a\)，即 \(z \leqslant a\)）组成的）非空。对 \(y \in A''\) 及 \(z \in A'\)，我们有 \(z \leqslant a \leqslant y\)，从而由假设 \(f(z) \leqslant f(y)\)。于是 \(\alpha', \alpha''\) 有限且 \(\alpha' \leqslant \alpha''\)。E 上延拓 f 的任何线性型 \(f_1\) 都由 \(f_1(a)\) 完全确定，且对所有 \(\lambda \in R\) 及所有 \(x \in V\)，我们有

\[
f _ {1} (x + \lambda a) = f (x) + \lambda f _ {1} (a).
\]

于是 \(f_{1}\) 是正的，当且仅当关系

\[
x \in \mathbf {V}, \quad \lambda \in \mathbf {R}, \quad x + \lambda a \geqslant 0\tag{2}
\]

蕴含

\[
f (x) + \lambda f _ {1} (a) \geqslant 0.\tag{3}
\]

由于 \(f(\mu x) = \mu f(x)\) 且关系 \(x \geqslant 0\) 与 \(\mu x \geqslant 0\)（对 \(\mu > 0\)）等价，只需证明 (2) 在特殊情形 \(\lambda = 0, \lambda = 1\) 与 \(\lambda = -1\) 下蕴含 (3)。对 \(\lambda = 0\)，(2) 蕴含 (3) 这一事实由 \(f\) 是正的这一假设得出。对 \(\lambda = 1\)，说 (2) 蕴含 (3) 意味着：对 \(-x \in A'\)，有 \(f_1(a) \geqslant f(-x)\)，即 \(f_1(a) \geqslant \alpha'\)；对 \(\lambda = -1\)，(2) 蕴含 (3) 意味着：对 \(x \in A''\)，有 \(f(x) \geqslant f_1(a)\)，即 \(f_1(a) \leqslant \alpha''\)。因此命题在这种情形下得证。

\paragraph{II. 一般情形。}

设 \(\mathfrak{F}\) 是由对 \((\mathrm{W}, g)\) 组成的集，其中 \(\mathrm{W}\) 是 \(\mathrm{E}\) 的包含 \(\mathrm{V}\) 的向量子空间，而 \(g\) 是 \(\mathrm{W}\) 上的正线性型且为 \(f\) 的延拓。我们给 \(\mathfrak{F}\) 排序：置 \((\mathrm{W}, g) \leqslant (\mathrm{W}', g')\)（若 \(\mathrm{W} \subset \mathrm{W}'\) 且 \(g'\) 是 \(g\) 的延拓）。显然 \(\mathfrak{F}\) 是归纳的，且由 S, III, § 2.4 的 th. 2，存在一个极大元素 \((\mathrm{W}_0, g_0)\)。设 \(\mathrm{W}_0 \neq \mathrm{E}\)。则存在向量 \(b \notin \mathrm{W}_0\)，且若 \(\mathrm{W}_1 = \mathrm{W}_0 + \mathbf{R}b\)，则上面的特殊情形表明存在 \(\mathrm{W}_1\) 上的正线性型且为 \(g_0\) 的延拓；这与 \((\mathrm{W}_0, g_0)\) 极大这一假设矛盾。于是 \(\mathrm{W}_0 = \mathrm{E}\)，命题的第一部分得证。当 \(a \in \mathrm{V}\) 时，第二个断言在 \(\alpha' = \alpha'' = f(a)\) 时显然成立；反之，若 \(a \notin \mathrm{V}\) 且置 \(\mathrm{V}_1 = \mathrm{V} + \mathbf{R}a\)，则第二个断言由证明中的特殊情形 I 得出。

推论. --- 在赋有相容预序结构的拓扑向量空间 E 中，设 P 是 E 中由 \(\geqslant 0\) 的元素组成的集。设 V 是 E 的一个向量子空间，至少包含 P 的一个内点 \(x_{0}\)。则 V 上每个正线性型都可以延拓为 E 上的正线性型。

由 prop. 1，只需证明对每个 \(x \in \mathbf{E}\)，存在 \(x' \in \mathbf{V}\) 使得 \(x' - x \in \mathbf{P}\)。现设 \(\mathbf{U}\) 是 \(\mathbf{E}\) 中 0 的一个邻域，使得 \(x_0 + \mathbf{U} \subset \mathbf{P}\)。则 \(x + x_0 + \mathbf{U} \subset x + \mathbf{P}\)，因此存在 \(\varepsilon\) 满足 \(0 < \varepsilon < 1\) 且点 \(y = x_0 + (1 - \varepsilon)x\) 属于 \(x + \mathbf{P}\)；于是每个形如 \(x + \lambda(y - x)\) 的点都属于 \(x + \mathbf{P}\)（对 \(\lambda > 0\)）。若取 \(\lambda = 1 / \varepsilon\)，则 \(x + \lambda(y - x) = \lambda x_0 \in \mathbf{V}\)，由此得出结论。

若不假设 V 包含 P 的一个内点，则推论的结论不一定成立，即使 E 是有限维的且 \(P \cap V\) 包含在 V 中为内点的点（II, 第 91 页, exerc. 25, b)）。
